% ============================================================================
% MASTER'S THESIS (GEDES) — SUBMISSION DRAFT
%
%   "Difference of Lunar Regolith Thermal Conductivity K_d
%    at the Apollo 15 and 17 Heat-Flow Boreholes"
%
% House rules (same as letter + guidebook):
%   * every number quoted here traces to results/*.json — never edit a value
%     by hand without re-deriving it from the pipeline;
%   * citations only from references.bib (symlinked to the letter's);
%   * figures via the figures/ symlink -> results/figures (single source);
%   * \todo{} markers flag passages only the author can write;
%   * amssymb must NOT be loaded (clashes with newtxmath on \Bbbk).
%
% Build:  latexmk -pdf thesis.tex
%
% Certified numbers baked in (2026-07-06 final audit; source letter.pdf):
%   K_d*(A15) = 4.60 (+0.86/-0.21) [4.18, 6.96]
%   K_d*(A17) = 7.08 (+0.53/-0.52) [6.16, 8.07]
%   contrast median +2.31, 95% CI [-0.12, 3.56], one-sided P_boot = 0.031
%   MCMC ordering P(A17 > A15) = 99.2-99.4% across prior widths
%   RMSE site-fit 1.00 / 0.40 K  vs global-3.4 fit 1.09 / 0.89 K
%   error-budget totals sigma = 1.88 / 3.88 mW m^-1 K^-1 (incl. sigma_epoch)
%   common-epoch: A17 -> 6.16 [5.92, 6.81], contrast +1.55, P = 0.04
%   speed: ~3000 lunations/brute solve; ~21x algorithmic, ~2500x compiled kernel
%          (reported as hardware-independent factors, NOT wall-clock times)
% ============================================================================
\documentclass[11pt,a4paper,oneside]{report}
% Float packing: let text share pages with floats instead of leaving
% near-empty float pages.
\renewcommand{\topfraction}{0.92}
\renewcommand{\bottomfraction}{0.6}
\renewcommand{\textfraction}{0.06}
\renewcommand{\floatpagefraction}{0.85}

\usepackage[T1]{fontenc}
\usepackage[utf8]{inputenc}
\usepackage{newtxtext,newtxmath}          % Times, matching the letter
\usepackage[margin=2.9cm]{geometry}
\usepackage{graphicx}
\usepackage{amsmath,bm}
\usepackage{siunitx}
\usepackage[round,authoryear]{natbib}
\usepackage{booktabs}
\usepackage{microtype}
\usepackage{xcolor}
\usepackage[colorlinks=true,linkcolor=blue!45!black,citecolor=blue!45!black,
            urlcolor=blue!45!black]{hyperref}
\usepackage{setspace}
\onehalfspacing

% ---- shared notation (canon of the letter) ---------------------------------
\newcommand{\Kd}{\ensuremath{K_d}}
\newcommand{\Kds}{\ensuremath{K_d^{*}}}
\newcommand{\Ks}{\ensuremath{K_s}}
\newcommand{\Qb}{\ensuremath{Q_b}}
\newcommand{\Tm}{\ensuremath{\langle T\rangle}}
\newcommand{\urect}{\ensuremath{u_{\mathrm{rect}}}}
\newcommand{\mWmK}{\si{\milli\watt\per\meter\per\kelvin}}
\newcommand{\mWm}{\si{\milli\watt\per\square\meter}}
\newcommand{\todo}[1]{\textcolor{red!70!black}{[TODO-#1]}}

% To show the institutional logo, drop the official Science Tokyo mark at
% figures/science_tokyo_logo.pdf (or .png); figures/ symlinks to the repo
% figures dir. It appears automatically once present; the build stays clean
% without it.
\title{%
  \IfFileExists{figures/science_tokyo_logo.pdf}
    {\includegraphics[height=1.7cm]{figures/science_tokyo_logo}\\[14pt]}
    {\IfFileExists{figures/science_tokyo_logo.png}
      {\includegraphics[height=1.7cm]{figures/science_tokyo_logo}\\[14pt]}{}}%
  \bfseries Difference of Lunar Regolith Thermal Conductivity
  $K_d$\\[4pt] at the Apollo 15 and 17 Heat-Flow Boreholes}
\author{Ramon III P.\ Gregorio \quad (Student No.\ 24M58378)\\[6pt]
  \normalsize Institute of Science Tokyo\\[6pt]
  \normalsize Supervisor: Yasuko Kasai\\
  \normalsize Co-supervisor: Richard Larsson}
\date{July 2026}

\begin{document}
\maketitle

% ============================================================================
\begin{abstract}
\noindent
Whether water ice can survive below the lunar surface, and how cold the
shallow ground stays, depend on the subsurface temperature, which no
instrument measures directly and every mission must model. The key
input to that model at depth is the deep thermal conductivity of the
regolith, \Kd, which sets how well the Moon's blanket of crushed rock
and dust insulates its interior. Every current model uses a single
Moon-wide value ($\Kd=3.4$~\mWmK; \citealp{hayne2017}) that has never
been checked against a subsurface measurement. This thesis measures
\Kd\ separately at the only two sites where the subsurface has been
sampled at meter depths, the restored Apollo 15 and 17 Heat Flow
Experiment (HFE) records \citep{nagihara2018}. Each sensor's multi-year
record is first reduced to one equilibrium temperature by an automatic
stability-window criterion; holding the standard model fixed
otherwise, the deep-sensor ($80$--$234$~cm) temperatures require
$\Kds(\mathrm{A15})=4.60^{+0.86}_{-0.21}\,[4.18,\,6.96]$ and
$\Kds(\mathrm{A17})=7.08^{+0.53}_{-0.52}\,[6.16,\,8.07]$~\mWmK\
(asymmetric $1\sigma$ statistical error, $95\%$ CI in brackets). The
total $1\sigma$ error budget is $\pm1.9$ and $\pm3.9$, dominated by
the uncertain basal heat flux and, at Apollo~17, the surface albedo. Both exceed the global value, and Apollo~17 is the more
conductive site: the inter-site contrast is marginal in size (its
$95\%$ interval spans zero) yet robust in direction across the full
published flux range. Underpinning this, the methodological
contribution is a flux-anchored equilibrium solver.
Once the subsurface settles into its repeating steady state, the
cycle-averaged heat flow is the same at every depth, so the deep
temperature profile can be rebuilt from a single anchor point rather
than simulated through thousands of day-night cycles, cutting each
solve's cost by about $2500\times$ and making a full uncertainty
analysis (bootstrap, degeneracy maps, Bayesian cross-checks) routine.
The retrieved temperature columns supply the meter-scale boundary
condition that subsurface-mapping missions such as the TSUKIMI
terahertz surveyor for lunar water and resources \citep{tsukimi2022}
cannot obtain from orbit, and the solver's speed opens the path from
these two calibration points toward Moon-wide maps of subsurface
temperature and of where water ice can persist.
\end{abstract}

\chapter*{Acknowledgments}
I would like to express my special thanks and deepest gratitude to my
advisor, Professor Yasuko Kasai, and my co-advisor, Richard Larsson,
for guiding me through this study. I also extend my thanks to my
colleagues at the Institute of Science Tokyo, Kasai Laboratory, and at
the National Institute of Information and Communications Technology
(NICT), for their valuable insights and helpful discussions.

I am deeply grateful to my scholarship provider, the Mindanao-Japan
Science and Technology Platform (MJ-STeP), and to Ateneo de Davao
University, for giving me the opportunity to pursue my graduate studies
in Japan. To my family, and to the friends I have met along the way,
thank you for the motivation and unwavering support you have given me
throughout this journey. And above all, to God, who made all of these
things possible.

\begin{quote}
\itshape\noindent
``For I know the plans I have for you,'' declares the Lord, ``plans to
prosper you and not to harm you, plans to give you hope and a future.''

\upshape\normalfont\raggedleft (Jeremiah 29:11)
\end{quote}

\tableofcontents
\listoffigures
\listoftables

% ============================================================================
\chapter*{Notation}
\addcontentsline{toc}{chapter}{Notation}

\noindent The symbols below keep the same meaning in every chapter and
match the companion letter and the code.

\medskip
\noindent\small\begin{tabular}{@{}llp{0.42\linewidth}@{}}
\toprule
Symbol & Value / units & Meaning \\
\midrule
\Kd            & \mWmK         & deep asymptote of the contact conductivity; the free parameter \\
\Kds           & \mWmK         & retrieved (best-fit) value of \Kd \\
\Ks            & $7.4\times10^{-4}$~\si{\watt\per\meter\per\kelvin} & surface conductivity \citep{hayne2017} \\
$H$            & $0.06$~m      & compaction e-folding depth \citep{hayne2017} \\
$\chi$         & $2.7$         & radiative conductivity coefficient \citep{hayne2017} \\
$T_{\rm ref}$  & $350$~K       & radiative normalization temperature \\
\Qb            & $21$ / $16$~\mWm & basal heat flux, A15 / A17 \citep{langseth1976} \\
$S_0$          & $1361$~\si{\watt\per\square\meter} & solar constant \citep{kopp2011} \\
$A$, $\varepsilon$ & $0.131$/$0.137$; $0.95$ & Bond albedo (A15/A17), emissivity \\
\Tm            & K             & cycle-mean (lunation-averaged) temperature \\
\urect         & \si{\watt\per\square\meter} & rectified flux, the mean-flux covariance correction \\
$T_{\rm eq}$   & K             & stability-window mean sensor temperature \\
$z_0$          & $0.55$~m      & anchor depth of the equilibrium solver \\
$\delta$       & $3$--$5$~cm   & diurnal thermal skin depth \\
$N$            & $7$ / $16$    & retained deep sensors, A15 / A17 \\
\bottomrule
\end{tabular}

\medskip
\noindent\begin{tabular}{@{}lp{0.66\linewidth}@{}}
\toprule
HFE & Heat Flow Experiment (the Apollo buried thermometers) \\
TG / TR & gradient-bridge / ring-bridge platinum sensor \\
RMSE & root-mean-square error over the retained deep sensors \\
AICc & small-sample-corrected Akaike information criterion \\
MCMC & Markov-chain Monte Carlo \\
GCP & Diviner Global Cumulative Product (surface temperatures) \\
RTM & radiative-transfer model \\
TSUKIMI & terahertz surveyor mission consuming the retrieved $T(z)$ \\
\bottomrule
\end{tabular}\normalsize

% ============================================================================
\chapter{Introduction}\label{ch:intro}

\section{A single Moon-wide conductivity}
The temperature of the lunar subsurface is not measured; it is
modeled. Every orbital instrument that senses below the optical
surface, from the Diviner radiometer
\citep{paige2010,williams2017diviner} to any future microwave or
terahertz sounder, receives radiation from depths no thermometer has
reached, and must invert what it sees through a forward model of the
temperature column $T(z)$. Below the top few centimeters, the single
most leveraged input to that column is the deep thermal conductivity
of the regolith (the Moon's surface layer of crushed rock and dust),
\Kd.

The standard source for \Kd\ is the global calibration of
\citet{hayne2017}, which fits a conductivity--density model to Diviner
nighttime cooling curves and arrives at $\Kd=3.4$~\mWmK, one value for
the entire Moon. The choice is reasonable for a global product, but
the diurnal wave Diviner sees, the temperature swing of the Moon's
day--night cycle, penetrates only a few skin depths, roughly $30$~cm.
At one to two meters the global value is an extrapolation, and it has
never been tested against an in-situ subsurface measurement, one made
in place on the Moon rather than inferred from orbit. The historical spread makes the
stakes concrete. Laboratory measurements of returned Apollo soils span
$1.2$--$3.5$~\mWmK\ \citep{cremers1971,hemingway1973,horai1981}; the
original in-situ HFE reductions gave $1$--$3$~\mWmK\
\citep{langseth1972,langseth1976}; later reanalyses revised these
upward \citep{keihm1984,grott2010}; orbital estimates reach
${\sim}7$~\mWmK\ \citep{vasavada2012}. Five decades of work disagree
by an order of magnitude (Fig.~\ref{fig:prior}).

\begin{figure}[htbp]
  \centering
  \includegraphics[width=0.8\linewidth]{figures/fig_prior_estimates.pdf}
  \caption[Five decades of meter-scale \Kd\ estimates]{Five decades of
  meter-scale conductivity estimates, separated by definition class.
  \textbf{(a)}~Effective (total) conductivities as sensors measure
  them, radiative term included (open markers): laboratory soils, the
  in-situ reductions, and orbital estimates. \textbf{(b)}~The contact
  asymptote \Kd, the Hayne-form model parameter (filled markers):
  the global fits and the per-site values retrieved in this work.
  This work appears in both panels, the two forms related by the
  radiative factor ${\sim}2$ at sensor temperatures; comparisons run
  within a panel, never across.}
  \label{fig:prior}
\end{figure}

The point of the figure is the spread. Laboratory, in-situ, and
orbital estimates scatter across nearly a decade with no convergence
in fifty years, and the single global fit that most models now adopt
sits inside that scatter without resolving it. The two panels keep
the two definition classes apart, the effective conductivity a sensor
measures and the contact parameter a model fits, which differ by a
factor of about two and are conflated surprisingly often in the
literature. The points labeled ``this work'' are what this thesis
adds: an in-situ value for each borehole site, in both classes, tied
to the only subsurface temperature data that exist.

\section{The Apollo Heat Flow Experiments as the sole in-situ constraint}
Two experiments constrain the meter-scale thermal state of the Moon in
situ: the Apollo 15 and Apollo 17 Heat Flow Experiments
\citep{langseth1972,langseth1976}. Each consists of fiberglass-epoxy
borestems carrying platinum-resistance sensor bridges to depths of
$1.4$ and $2.3$~m, which returned five and a half years (1971--1977)
of subsurface temperatures, restored to modern usability by
\citet{nagihara2018}. One further in-situ measurement now exists: the
ChaSTE probe on the Chandrayaan-3 lander recorded a temperature
profile over the top ${\sim}10$~cm at a high-latitude site
\citep{mathew2025chaste}, valuable in its own right but still well
within the diurnal skin. The two Apollo boreholes remain the only
meter-scale subsurface profiles. Figure~\ref{fig:context} locates the two sites
on the nearside: Apollo~15 at the Hadley--Apennine mare/highland
margin ($26.1^{\circ}$N, $3.6^{\circ}$E) and Apollo~17 in the
Taurus--Littrow valley ($20.2^{\circ}$N, $30.8^{\circ}$E).

\begin{figure}[htbp]
  \centering
  \includegraphics[width=\linewidth]{figures/fig_context_globes.pdf}
  \caption[The two Apollo HFE sites on the lunar nearside]{The two
  Apollo HFE landing sites on the lunar nearside, shown as three data
  layers of the same Earth-facing disk. \textbf{(a)}~Clementine albedo
  mosaic; \textbf{(b)}~modeled diurnal-mean surface temperature from
  the solver of Chapter~\ref{ch:model}; \textbf{(c)}~topography from
  LOLA (the Lunar Orbiter Laser Altimeter). Both boreholes are marked
  on every globe.}
  \label{fig:context}
\end{figure}

Read across the three layers, the two sites are not twins. Apollo~15
sits on the boundary between bright highlands and dark mare, on
regolith mixed from both; Apollo~17 lies in a low, dark basaltic
valley closed in by highland massifs
\citep{korotev2005,wieczorek2006,crawford2015}. The two columns have
different provenance and compaction histories. That is the first
reason to doubt that one deep conductivity should fit both.

This thesis treats the two boreholes as what they are, the only
meter-scale ground truth available, and puts the narrowest useful
question to them: with the community-standard conductivity model held
fixed in every other respect, what deep conductivity each borehole
requires. Section~\ref{sec:intro-tsukimi} explains why the answer
matters, and the following section states it as three precise
research questions.

\section{Relevance to subsurface radiative-transfer retrievals}\label{sec:intro-tsukimi}
The term \emph{meter-scale regolith} is used throughout for the layer
constrained by the deep HFE sensors, $80$--$234$~cm. Below the diurnal
skin, the shallow layer that the day--night temperature swing actually
reaches, the modeled column is controlled by \Kd\ through two channels,
the rectified lunation mean and the steady geothermal gradient
$\Qb/\Kd$, so the $T(z)$ ingested by any subsurface
radiative-transfer (RTM) retrieval inherits its \Kd\ dependence. A
factor-of-two error in \Kd\ propagates through the meter-scale
temperature profile into every product built on it.

A concrete consumer is the upcoming TSUKIMI mission (Lunar Terahertz
SUrveyor for KIlometer-scale MappIng; \citealp{tsukimi2022}), a
NICT/JAXA-led terahertz surveyor for subsurface water and resource
mapping. Terahertz emission originates below the diurnal skin, so the
TSUKIMI forward model needs precisely the sub-skin temperature column
that only an in-situ-anchored \Kd\ can supply. The same requirement
applies to any future microwave sounder.

\section{Research questions and contributions}
Three questions organize the work.
\begin{enumerate}
  \item Given the restored HFE records, what deep conductivity does
    each borehole require under the standard \citet{hayne2017} model,
    and with what uncertainty?
  \item Do the two sites require the \emph{same} value, as every
    global model assumes?
  \item Can the answer be made computationally cheap enough that the
    full uncertainty analysis (bootstrap, degeneracy maps, Bayesian
    cross-checks) is routine rather than heroic?
\end{enumerate}

The thesis answers them with three contributions.
\begin{enumerate}
  \item \textbf{A method.} The flux-anchored equilibrium solver
    (Chapter~\ref{ch:model}) converts the periodic steady-state
    problem from thousands of time-stepped lunations (a lunation is
    one lunar day--night cycle, about $29.5$ Earth days) into a
    skin-only time-step plus a one-variable closure ODE: a measured $21\times$
    acceleration per solve at the algorithmic level, ${\sim}2500\times$
    with the compiled kernel, at certified-equal answers.
  \item \textbf{A measurement.} Per-site deep conductivities
    $\Kds(\mathrm{A15})=4.60$ and $\Kds(\mathrm{A17})=7.08$~\mWmK\
    with a bootstrap-based statistical uncertainty and a full
    systematic budget (Chapters~\ref{ch:retrieval}
    and~\ref{ch:results}).
  \item \textbf{A robustness case.} A direct map of the $\Kd$--$\Qb$
    degeneracy, one full re-retrieval per grid point, showing that the
    two sites respond to basal-flux revisions in opposite directions.
    This makes the \emph{sign} of the inter-site contrast robust
    across the entire published flux envelope even though its
    magnitude is marginal (Section~\ref{sec:robustness}).
\end{enumerate}

\section{Thesis structure}
Chapter~\ref{ch:data} describes the data: the instrument, the restored
record, and the reduction of six years of temperatures to one
equilibrium estimate per sensor. Chapter~\ref{ch:model} builds the
thermal model and the flux-anchored solver.
Chapter~\ref{ch:retrieval} defines the retrieval and its uncertainty
treatment; Chapter~\ref{ch:results} reports the results and the
robustness program. Chapter~\ref{ch:discussion} interprets the
contrast, states the caveats, and closes with implications for future
missions. Appendix~\ref{app:derivation} derives the closure ODE in
full; Appendix~\ref{app:numerics} records the numerical scheme and its
certification, including the audit trail of the three input errors
found and fixed on the way to the final numbers;
Appendix~\ref{app:repro} gives reproduction instructions. All code
and data references live in the appendices.

Two reading paths are supported. A reader interested in the
measurement can go from Chapter~\ref{ch:data} directly to
Chapters~\ref{ch:retrieval} and~\ref{ch:results}, taking the solver
on trust from the certification summary at the end of
Chapter~\ref{ch:model}. A reader interested in the method can treat
Chapter~\ref{ch:model} and Appendices~\ref{app:derivation}
and~\ref{app:numerics} as a self-contained unit: nothing in the
flux-anchored construction depends on which borehole it is applied
to, and the certification argument (convergence plateaus in the
reported quantity, an independent implementation, a time-step ladder)
transfers to any periodically forced conduction problem.

Figure~\ref{fig:pipeline} condenses the whole program into one chart.
The spine runs from the restored record to the headline numbers; the
loop on the right is the sweep, one fast equilibrium solve per
\Kd\ candidate, which is where the fast solver earns its place. Each
box names the chapter that owns it, so the figure doubles as a map of
the thesis.

\begin{figure}[htbp]
  \centering
  \includegraphics[width=0.82\linewidth]{figures/fig_pipeline_overview.pdf}
  \caption[The whole pipeline in one chart]{The retrieval pipeline
  from the restored HFE record to the headline result, with the
  owning chapter tagged in each box. The loop on the right is the
  \Kd\ sweep: one flux-anchored equilibrium solve per candidate
  value. Teal boxes are the computational core, coral the uncertainty
  machinery, green the result.}
  \label{fig:pipeline}
\end{figure}

% ============================================================================
\chapter{The Apollo Heat Flow Experiment dataset}\label{ch:data}

\section{Instrument design and measurement geometry}
Each HFE site carries two probes about $9$~m apart. Each probe is a
fiberglass-epoxy borestem holding gradient-bridge (TG) and
ring-bridge (TR) platinum-resistance sensors, with thermocouples (TC)
in the shallow cable connecting it to the central ALSEP station
(Fig.~\ref{fig:probes}). A gradient
bridge reads the temperature difference across a known vertical
baseline of roughly $50$~cm, which is what makes the experiment a
\emph{heat-flow} experiment: with a conductivity model, the flux
$q=-K\,\partial T/\partial z$ follows directly. Apollo~15 (Hadley
Rille) reaches $1.39$~m; Apollo~17 (Taurus--Littrow) reaches
$2.34$~m. Sensor placement depths carry a ${\sim}2.5$~cm uncertainty
\citep{nagihara2018}, which the bootstrap of
Chapter~\ref{ch:retrieval} propagates explicitly.

Two features of the design matter for what follows.
First, the borestem is fiberglass-epoxy, a poor conductor chosen so
the tube perturbs the surrounding regolith as little as possible;
even so, it conducts the surface temperature wave downward well enough
to disturb the shallowest sensors, which is the reason the retrieval
later excludes everything above $80$~cm (\S\ref{sec:deep-cut}).
Second, the gradient bridge does not read an absolute temperature but
the difference $\Delta T$ across its baseline, so the measured
quantity is already a temperature gradient; a conductivity is all that
stands between it and a heat flow.
% AUTHOR TODO (optional, currently omitted): a real deployed-HFE
% photograph (NASA public-domain) showing the borestem and its cable to
% the ALSEP station could be inserted here as a figure.

\begin{figure}[htbp]
  \centering
  \includegraphics[width=\linewidth]{figures/fig_intro_probe.pdf}
  \caption[Apollo HFE measurement geometry and the published $K(z)$
  forms]{Apollo HFE measurement geometry and the two published $K(z)$
  forms tested. \textbf{(a)}~Borestem emplacements at both sites to a
  shared depth scale: markers are the real sensor depths, filled where
  a sensor enters the retrieval ($z\ge80$~cm) and open in the
  borestem-disturbed zone ($z<80$~cm, coral); the two probe strings
  per site (P1 left, P2 right) are offset. Apollo~17 reaches $2.3$~m
  (16 retained sensors), Apollo~15 $1.4$~m (7).
  \textbf{(b)}~Conductivity profiles $K(z)$ at $T=250$~K:
  \citet{hayne2017} at the published global $\Kd=3.4$~\mWmK\ and the
  \citet{martinez2021} $K(T,\rho)$ form, both introduced in
  Chapter~\ref{ch:model}.}
  \label{fig:probes}
\end{figure}

The two panels of Fig.~\ref{fig:probes} already contain the shape of
the whole problem. Panel~(a) shows where the information lives: a
handful of sensors below $80$~cm, more numerous and deeper at
Apollo~17. Panel~(b) shows what is being tested: two published
conductivity profiles that agree near the surface and settle onto
different deep asymptotes, exactly the region the deep sensors probe.

\section{The restored 1971--1977 record and its disturbances}
The HFE returned data continuously from deployment (July 1971 at
Apollo~15, December 1972 at Apollo~17) until the ALSEP (Apollo Lunar
Surface Experiments Package) stations were switched off in September
1977. For decades only the 1971--1974
portion was archived; \citet{nagihara2018} recovered the remainder
from the original tapes, and their restored record is the data source
of this thesis. The record is not quiet. Drilling and emplacement
heated each borehole, the early months show the decay of that
transient, and documented operational events (ALSEP outages, station
commands, eclipse responses) interrupt the series
(Figs.~\ref{fig:timelineA15} and~\ref{fig:timelineA17}).

\begin{figure}[t]
  \centering
  \includegraphics[width=0.92\linewidth]{figures/fig_apollo_timeline_a15.pdf}
  \caption[The restored Apollo 15 record and stability windows]{The
  restored 1971--1977 Apollo 15 record, both probes. Shaded vertical
  bands mark documented operational disturbances (drilling transient,
  ALSEP outages, station events); the stability windows actually used
  are highlighted, and the dashed line is the least-squares slope of
  each probe's deepest sensor ($+0.065$ and $+0.078$~K~yr$^{-1}$, both
  below the $0.08$~K~yr$^{-1}$ threshold). Gray rows are
  borestem-excluded sensors.}
  \label{fig:timelineA15}
\end{figure}

\begin{figure}[t]
  \centering
  \includegraphics[width=0.92\linewidth]{figures/fig_apollo_timeline_a17.pdf}
  \caption[The restored Apollo 17 record and stability windows]{The
  restored Apollo 17 record, both probes, in the same layout as
  Fig.~\ref{fig:timelineA15}. The deepest-sensor slopes are $+0.060$
  and $+0.011$~K~yr$^{-1}$; the deeper, colder-start Apollo 17
  borehole relaxes more slowly, which is why its stability windows
  open later in the mission.}
  \label{fig:timelineA17}
\end{figure}

\section{Reduction to per-sensor equilibrium temperatures}\label{sec:stable-window}
The retrieval does not fit the raw time series. It fits one number per
sensor: the equilibrium temperature $T_{\rm eq}$, the mean over a
late-mission window in which the record is statistically flat. The
window is selected per sensor by a fixed rule. Thirteen equally spaced
candidate start points spanning $55$--$80\%$ of the record are each
given an ordinary least-squares fit through everything from that point
to the end (candidates keeping less than $20\%$ of the record are
discarded). The accepted window is the earliest candidate whose trend
satisfies $|d\langle T\rangle/dt|<0.08$~K~yr$^{-1}$, the longest tail
consistent with thermal equilibrium at that stringency. The threshold
is roughly $4.6$ times stricter than the $10^{-3}$~K~day$^{-1}$
criterion sometimes used in HFE work, chosen to suppress bias from
residual post-drilling drift. Where no candidate passes, the final
$25\%$ of the record is used as a fallback. The strict criterion is
met by 3 of the 7 retained Apollo~15 sensors and 6 of the 16 at
Apollo~17; the rest take the fallback window and carry residual
trailing slopes of $0.11$--$0.47$~K~yr$^{-1}$, a systematic that
Chapter~\ref{ch:results} measures directly
(\S\ref{sec:commonepoch}).

Figure~\ref{fig:windowanat} shows the rule operating on two real
Apollo~15 sensors, one for each of its branches. In panel~(a) the
deepest sensor's record relaxes from the drilling spike, and by the
first candidate point the remaining trend is already
$+0.065$~K~yr$^{-1}$, under the threshold, so the longest candidate
window is accepted outright. In panel~(b) a shallower sensor is still
visibly warming: all thirteen candidate trends exceed the threshold
(the first fits at $+0.23$~K~yr$^{-1}$), so the rule falls back to the
final quarter of the record and the residual $+0.22$~K~yr$^{-1}$ slope
is carried as a data systematic rather than hidden. The two panels are
the entire estimator: every $T_{\rm eq}$ in Table~\ref{tab:sensors}
came from one of these two paths.

\begin{figure}[htbp]
  \centering
  \includegraphics[width=0.9\linewidth]{figures/fig_window_anatomy.pdf}
  \caption[The stability-window rule on two real sensors]{The
  stability-window rule operating on two real Apollo~15 records.
  Vertical ticks: the 13 candidate start points spanning
  $55$--$80\%$ of each record. \textbf{(a)}~The deepest sensor
  (TG12B, $139$~cm): the first candidate's trend is already below the
  $0.08$~K~yr$^{-1}$ threshold, so the longest window is accepted
  (green). \textbf{(b)}~A still-relaxing sensor (TR12A, $101$~cm):
  every candidate exceeds the threshold, so the rule takes the
  final-quarter fallback (coral) and its residual slope enters the
  error budget.}
  \label{fig:windowanat}
\end{figure}

The choice of threshold barely matters. Sweeping it across
$0.04$--$0.16$~K~yr$^{-1}$, a factor of four around the adopted value,
moves the retrieved \Kds\ by at most $0.02$~\mWmK\ at Apollo~15 and
$0.05$~\mWmK\ at Apollo~17, an order of magnitude below the bootstrap
uncertainty. The window rule is a discretionary choice sitting on a
flat response.

One caution applies to everything downstream. These $T_{\rm eq}$
values are estimates of the undisturbed periodic-equilibrium profile,
not proof that every trace reached exact steady state. The boreholes
were heated by drilling and relaxed toward the surrounding regolith
for years; the model is therefore compared to
equilibrium-temperature estimates, and the residual drift of the
fallback-window sensors is carried as a data systematic rather than
ignored.

\section{Selection of the meter-scale sensor set}\label{sec:deep-cut}
Sensors above $80$~cm are excluded from the retrieval, for two
independent reasons.

The first is instrumental. The borestem conducts the surface wave
downward and thermally short-circuits the upper regolith, the
long-recognized borestem problem of the original reductions
\citep{langseth1976,keihm1984,grott2010}. The diurnal amplitude
diagnostic of Fig.~\ref{fig:amplitude} makes the contamination visible
without reference to any residual: above $80$~cm the observed
amplitudes exceed the modeled regolith attenuation by an order of
magnitude or more, the signature of heat carried down the stem.

The second is physical. The diurnal wave attenuates exponentially
with depth, and by ${\sim}80$~cm its amplitude is far below sensor
noise. Deep sensors therefore read the cycle-mean state that \Kd\
controls, uncontaminated by the shallow-skin dynamics that \Ks\
governs.

\begin{figure}[htbp]
  \centering
  \includegraphics[width=0.82\linewidth]{figures/fig_amplitude_vs_depth.pdf}
  \caption[Diurnal amplitude vs.\ depth: the borestem
  diagnostic]{Observed diurnal amplitude (estimated as
  $\sqrt{2}\times$ the within-window standard deviation) against
  sensor depth at both sites, on a logarithmic amplitude axis, with
  the modeled regolith attenuation curves for comparison. Open
  markers: borestem-zone sensors ($z<80$~cm); filled markers: the
  retained deep set. Above the shaded borestem zone the observed
  amplitudes sit an order of magnitude or more above the attenuation
  curves, the heat-short signature of the conducting stem; below it
  they collapse onto the $0.05$--$0.2$~K digitization floor of the
  restored record.}
  \label{fig:amplitude}
\end{figure}

Figure~\ref{fig:amplitude} shows both effects in one frame. Reading
either site's markers downward, the observed amplitude falls by more
than two decades from the surface to the borestem line, but far more
slowly than the modeled regolith attenuation demands: the shallow
markers ride an order of magnitude or more above the curves, because
the fiberglass stem carries the surface wave down past them. At the
$80$~cm line the observed points meet the digitization floor and stop
falling, and every deeper sensor sits on that floor. The cut is placed
where the contamination signal ends, and the amplitude information
ends with it: below the line, both published models predict amplitudes
more than five orders of magnitude under the floor, so deep amplitudes
carry no diffusivity information and the diagnostic sorts clean
sensors from contaminated ones without touching the quantity being
retrieved.

The cut leaves $N=7$ deep sensors at Apollo~15 ($80$--$139$~cm) and
$N=16$ at Apollo~17 ($130$--$234$~cm). The retrieval is insensitive to
the exact depth of the cut: sweeping it over $70$--$90$~cm holds
$\Kds(\mathrm{A15})$ at $4.6$--$4.8$ and $\Kds(\mathrm{A17})$ at
$7.1$~\mWmK, a contrast of $2.2$--$2.5$~\mWmK, unchanged within
rounding. Only a $60$~cm cut deviates ($\Kds(\mathrm{A17})=8.3$,
contrast $3.6$), because two borestem-affected sensors re-enter the
Apollo~17 set and inflate the gradient, exactly the contamination the
amplitude diagnostic exists to exclude. The outlier confirms the
diagnostic rather than undermining it.

\section{Equilibrium temperature profiles}\label{sec:eqprofiles}
Figure~\ref{fig:meanT} shows the observable this thesis fits: a set of
late-window equilibrium temperatures against depth, not a movie, and
Table~\ref{tab:sensors} lists every retained value behind it. Three
features of the figure carry the physics. First, the two sites
separate cleanly: the Apollo~15 deep sensors cluster near $253$~K
while the Apollo~17 set runs near $256$~K, about $3$~K warmer through
the meter-scale column. Second, within each site the filled deep
points climb along a gradient that the excluded shallow sensors do not
share; that deep gradient, not the absolute temperature, is what \Kd\
sets and what the retrieval reads. Third, the scatter about each
site's trend is not purely instrumental: interleaved sensors from the
two probe strings differ by up to ${\sim}0.6$~K at the same depth
(the two $140$~cm Apollo~17 sensors read $255.7$ and $256.2$~K), a
probe-to-probe spread the bootstrap of Chapter~\ref{ch:retrieval}
treats as sensor-set scatter.

\begin{figure}[htbp]
  \centering
  \includegraphics[width=0.8\linewidth]{figures/fig_apollo_mean_T_profile.pdf}
  \caption[Mean temperature profiles at both sites]{Stable-window mean
  temperatures against depth at both sites, with the borestem
  exclusion zone shaded and the smooth model profiles for context.
  Filled markers: the deep TG/TR sensors that enter the retrieval
  (charcoal TG, gray TR); open markers: shallow borestem-zone sensors,
  excluded. Apollo~17 runs about $3$~K warmer than Apollo~15 at every
  retained depth, and each site's filled points define the deep
  gradient the retrieval fits.}
  \label{fig:meanT}
\end{figure}

\begin{table}[htbp]
  \centering\footnotesize
  \setlength{\tabcolsep}{4.5pt}
  \caption[The retained deep-sensor set]{The retained deep-sensor set:
  every temperature the retrieval fits, extracted by the pipeline of
  Appendix~\ref{app:repro}. ``flat'' windows satisfy the
  $0.08$~K~yr$^{-1}$ criterion of \S\ref{sec:stable-window};
  ``tail'' sensors use the final-quarter fallback window, and their
  residual trailing slope is listed.}
  \label{tab:sensors}
  \begin{tabular}{llrrrlr@{\qquad}llrrrlr}
    \toprule
    \multicolumn{7}{c}{Apollo 15 ($N=7$)} &
    \multicolumn{7}{c}{Apollo 17 ($N=16$)} \\
    \cmidrule(r{2em}){1-7}\cmidrule{8-14}
    P & type & $z$ & $T_{\rm eq}$ & $\sigma$ & win & slope &
    P & type & $z$ & $T_{\rm eq}$ & $\sigma$ & win & slope \\
      &      & cm  & K            & K        &     & K/yr  &
      &      & cm  & K            & K        &     & K/yr  \\
    \midrule
    1 & TG & 84  & 253.20 & 0.05 & tail & $+0.13$ & 1 & TG & 130 & 256.48 & 0.04 & tail & $+0.32$ \\
    2 & TR & 87  & 250.88 & 0.09 & tail & $+0.19$ & 2 & TG & 131 & 256.59 & 0.04 & tail & $+0.22$ \\
    1 & TG & 91  & 253.15 & 0.04 & flat & ---     & 1 & TR & 140 & 255.67 & 0.08 & tail & $+0.47$ \\
    2 & TG & 97  & 251.16 & 0.05 & flat & ---     & 2 & TR & 140 & 256.24 & 0.04 & tail & $+0.23$ \\
    1 & TR & 101 & 253.14 & 0.07 & tail & $+0.22$ & 1 & TR & 167 & 255.88 & 0.06 & tail & $+0.33$ \\
    1 & TR & 129 & 253.36 & 0.06 & tail & $+0.19$ & 2 & TR & 169 & 256.38 & 0.03 & tail & $+0.15$ \\
    1 & TG & 139 & 253.53 & 0.03 & flat & ---     & 1 & TG & 177 & 256.60 & 0.03 & tail & $+0.22$ \\
      &    &     &        &      &      &         & 2 & TG & 178 & 256.74 & 0.03 & tail & $+0.11$ \\
      &    &     &        &      &      &         & 1 & TG & 185 & 256.06 & 0.05 & tail & $+0.22$ \\
      &    &     &        &      &      &         & 2 & TG & 186 & 256.52 & 0.03 & flat & ---     \\
      &    &     &        &      &      &         & 1 & TR & 195 & 256.13 & 0.04 & tail & $+0.19$ \\
      &    &     &        &      &      &         & 2 & TR & 196 & 256.55 & 0.02 & flat & ---     \\
      &    &     &        &      &      &         & 1 & TR & 223 & 256.38 & 0.03 & flat & ---     \\
      &    &     &        &      &      &         & 2 & TR & 224 & 256.74 & 0.02 & flat & ---     \\
      &    &     &        &      &      &         & 1 & TG & 233 & 256.48 & 0.03 & flat & ---     \\
      &    &     &        &      &      &         & 2 & TG & 234 & 256.84 & 0.03 & flat & ---     \\
    \bottomrule
  \end{tabular}
\end{table}

The table also previews two facts used later. The ``flat'' windows
concentrate at depth: every Apollo~17 sensor that passes the strict
criterion sits at $186$~cm or deeper, while the shallower rows ride
late-relaxing fallback tails; that depth-aligned pattern is the root
of the epoch mixing measured in \S\ref{sec:commonepoch}. And the
per-sensor scatter $\sigma$ of $0.02$--$0.09$~K is an order of
magnitude smaller than the $0.4$--$1$~K model residuals of
Chapter~\ref{ch:results}, which is why model adequacy, not thermometer
noise, dominates the fit statistics (\S\ref{sec:aicc-method}).

\section{Basal heat-flux boundary condition}\label{sec:qb-data}
The lower boundary condition is the basal heat flux \Qb, the steady
trickle of heat flowing up from the Moon's interior, at the revised
borehole values of \citet{langseth1976}: $21$~\mWm\ at Apollo~15 and
$16$~\mWm\ at Apollo~17. These values are debated: \citet{saito2007}
argue for fluxes up to a third lower, and \citet{nagihara2018} note
limits in the underlying reduction. The debate matters because the
deep gradient constrains $\Qb/\Kd$, so any flux revision moves \Kds.
Section~\ref{sec:robustness} therefore treats \Qb\ as the dominant,
explicitly mapped systematic, re-retrieving \Kds\ over its full
published envelope
($\Qb^{\mathrm{A15}}\in[14,25]$, $\Qb^{\mathrm{A17}}\in[10,18]$~\mWm).

% ============================================================================
\chapter{The thermal model and the flux-anchored
  solver}\label{ch:model}\label{ch:method}

\section{Governing equation and thermophysical properties}
Heat moves through regolith by two mechanisms. Contact conduction
carries phonons through the network of grain-to-grain contacts, so it
grows with compaction and hence with depth. Radiative exchange carries
infrared photons across the pore spaces between grains, so it grows
steeply with temperature, as $T^3$. The \citet{hayne2017} model wraps
both into one field,
\begin{equation}
  K(T,z) = \underbrace{\left[\Kd - (\Kd-\Ks)\,
      e^{-z/H}\right]}_{\text{contact conduction } K_c(z)}
    \underbrace{\left[1+\chi\left(\frac{T}{T_{\rm ref}}\right)^{3}
      \right]}_{\text{radiative factor}},
  \label{eq:conductivity}
\end{equation}
and the subsurface temperature obeys the one-dimensional heat
equation
\begin{equation}
  \rho(z)\, c_p(T)\, \frac{\partial T}{\partial t}
  = \frac{\partial}{\partial z}\!\left[K(T,z)\,
    \frac{\partial T}{\partial z}\right],
  \label{eq:heat}
\end{equation}
with the density $\rho(z)$ relaxing from a surface to a deep value on
the same compaction scale $H$ and the specific heat $c_p(T)$ from the
\citet{hemingway1973} polynomial. Table~\ref{tab:params} lists every
fixed input with its source. \Kd, the deep asymptote of the contact
term, is the single free parameter of this work. Every other constant
is held at its published value, deliberately, so that the retrieved
\Kds\ is interpretable \emph{within} the standard model rather than as
a new model.

\begin{table}[htbp]
  \centering\small
  \caption[Fixed inputs to the thermal model]{Fixed inputs to the
  thermal model. Site-specific rows are listed A15\,/\,A17. The
  \citet{martinez2021} coefficients enter only the comparison model of
  Chapter~\ref{ch:results}.}
  \label{tab:params}
  \begin{tabular}{@{}llll@{}}
    \toprule
    Parameter & Symbol & Value & Source \\
    \midrule
    Surface conductivity & \Ks & $7.4\times10^{-4}$~\si{\watt\per\meter\per\kelvin} & \citet{hayne2017} \\
    Deep conductivity (global) & \Kd & $3.4\times10^{-3}$~\si{\watt\per\meter\per\kelvin} & \citet{hayne2017} \\
    Compaction scale & $H$ & $0.06$~m & \citet{hayne2017} \\
    Radiative coefficient & $\chi$ & $2.7$ & \citet{hayne2017} \\
    Radiative normalization & $T_{\rm ref}$ & $350$~K & \citet{hayne2017} \\
    Surface / deep density & $\rho_s$, $\rho_d$ & $1100$, $1800$~kg\,m$^{-3}$ & \citet{hayne2017} \\
    Specific heat & $c_p(T)$ & 4th-order polynomial & \citet{hemingway1973} \\
    Latitude & $\phi$ & $26.1^{\circ}$\,/\,$20.2^{\circ}$\,N & ALSEP gazetteer \\
    Bond albedo & $A$ & $0.131$\,/\,$0.137$ & \citet{vasavada2012} \\
    Emissivity & $\varepsilon$ & $0.95$ & \citet{bandfield2015} \\
    Basal heat flux & \Qb & $21$\,/\,$16$~\mWm & \citet{langseth1976} \\
    Solar constant & $S_0$ & $1361$~\si{\watt\per\square\meter} & \citet{kopp2011} \\
    \midrule
    Mart\'inez contact coeffs. & $A_1$, $A_2$ & $5.082\times10^{-6}$, $-5.10\times10^{-3}$ & \citet{martinez2021} \\
    Mart\'inez radiative coeffs. & $B_1$, $B_2$ & $2.002\times10^{-13}$, $-1.953\times10^{-10}$ & \citet{martinez2021} \\
    \bottomrule
  \end{tabular}
\end{table}

Figure~\ref{fig:kTz} shows the $K(T,z)$ field under the retrieved
conductivities, zoomed into the top $40$~cm where all the structure
lives. Two features stand out: the steep vertical rise as the contact
conductivity climbs from its tiny surface value to the deep asymptote
within a few $H$, and the horizontal warm tilt of the radiative
$\chi(T/T_{\rm ref})^3$ factor. Below this window the field is uniform
at \Kd, which is the flat state the deep HFE sensors sample.

\begin{figure}[htbp]
  \centering
  \includegraphics[width=0.9\linewidth]{figures/fig_book_kTz.pdf}
  \caption[The $K(T,z)$ field under the retrieved conductivities]{The
  \citet{hayne2017} conductivity field $K(T,z)$ evaluated with the
  retrieved per-site \Kds, zoomed to the top $40$~cm. The steep
  near-surface rise is the compaction term on the scale $H$; the
  temperature dependence is the radiative $\chi(T/T_{\rm ref})^3$
  factor. Below this window the field is uniform at \Kd.}
  \label{fig:kTz}
\end{figure}

\section{Boundary conditions}\label{sec:bc}
The Moon has no atmosphere, so the surface exchanges energy by
radiation and conduction alone. The solar irradiance on the local
horizontal is
\begin{equation}
  F_\odot(t) = \max\bigl(0,\ S_0\cos\theta_\odot(t)\bigr),
  \label{eq:insolation}
\end{equation}
with the daytime zenith angle approximated by the mean synodic cycle
at the site latitude,
$\cos\theta_\odot(t)=\cos\phi\,\cos(2\pi t/P)$, $P=29.531$~d. At the
surface, absorbed insolation, thermal emission, and conduction into
the ground must balance,
\begin{equation}
  (1-A)\,F_\odot(t)\;-\;\varepsilon\sigma T_s^4
  \;-\;\bigl(-K\,\partial_z T\bigr)_{z=0} \;=\; 0,
  \label{eq:bc-top}
\end{equation}
a nonlinear equation for the skin temperature $T_s$ solved by Newton
iteration at every time step; here $A$ is the Bond albedo, the
fraction of sunlight the surface reflects. At the base ($z_{\max}=5$~m)
the geothermal flux enters as a Neumann condition, one that fixes the
heat flow rather than the temperature:
$K\,\partial_z T|_{z_{\max}}=\Qb$, imposed through a ghost cell.
Figure~\ref{fig:bc} sketches both boundaries. A fixed-temperature or
zero-flux base would be wrong: the deep profile keeps rising at the
geothermal gradient $\Qb/K$, roughly $2.3$~K~m$^{-1}$ at Apollo~15,
a physical feature of the steady state rather than a numerical
artifact.

\begin{figure}[t]
  \centering
  \includegraphics[width=0.9\linewidth]{figures/fig_boundary_conditions.pdf}
  \caption[How the Sun and the interior bound the column]{How the Sun
  and the lunar interior bound the model column. At the surface, the
  absorbed insolation $(1-A)F_\odot$ balances thermal emission
  $\varepsilon\sigma T_s^4$ and conduction into the ground; the
  balance is closed for $T_s$ by Newton iteration at every time step.
  At the base, the geothermal flux enters as the Neumann condition
  $K\,\partial_z T=\Qb$ ($21$ and $16$~\mWm\ at Apollo 15 and 17),
  which sets the deep gradient the retrieval reads.}
  \label{fig:bc}
\end{figure}

The nonlinearity of Eq.~\eqref{eq:bc-top} is mild but real: the
emission term grows as $T_s^4$, so the balance cannot be solved in
closed form. Figure~\ref{fig:newton} shows the residual at one noon
instant of the converged Apollo~17 cycle. The curve crosses zero once,
at $T_s=377$~K, and Newton iteration walks to that root in three steps
even from a deliberately distant starting guess; in production each
time step starts from the previous step's $T_s$ and closes in a single
step. This small nonlinear solve runs inside every one of the
half-hour steps of the march described next.

\begin{figure}[htbp]
  \centering
  \includegraphics[width=0.62\linewidth]{figures/fig_newton_surface.pdf}
  \caption[Newton iteration on the surface balance]{The surface energy
  balance residual against trial skin temperature at one noon instant
  of the converged Apollo~17 cycle. The $T_s^4$ emission term steepens
  the curve; Newton iteration (coral tangent steps, started from a
  deliberately distant guess to display the iteration) reaches the
  root at $T_s=377$~K in three steps. Production restarts from the
  previous step's $T_s$ and lands in one.}
  \label{fig:newton}
\end{figure}

\section{Spatial and temporal discretization}\label{sec:discretization}
The continuous problem of Eqs.~\eqref{eq:heat}--\eqref{eq:bc-top} is
solved on a geometric finite-volume grid, and one design decision
carries the whole discretization: resolution goes where the curvature
is. The ${\sim}280$~K diurnal swing lives in the top few centimeters,
so the first cell is $\Delta z_0=2$~mm thick and each cell below grows
by $8\%$, reaching the $5$~m base in 69 cells. Temperatures live at
cell centers and fluxes cross cell faces, so integrating the heat
equation over a cell conserves energy exactly; the face conductivity
is the harmonic mean of the two neighboring cells, because heat
crossing a face passes through two half-cells in series and thermal
resistances add. Where $K$ climbs steeply through the compaction zone
an arithmetic mean would let a high-$K$ cell short-circuit its
low-$K$ neighbor.

Figure~\ref{fig:numerics} lays the scheme out. Panel~(a) plots the
real grid, cell width against depth: a straight line on logarithmic
axes, which is what geometric growth looks like, running from
millimeter cells under the surface to decimeter cells in the smooth
deep column, with the anchor and the sensor band sitting comfortably
in the well-resolved middle. Panel~(b) is the stencil: each cell
exchanges flux with its two neighbors through harmonic-mean faces, so
each interior equation couples three unknowns. Panel~(c) is the
consequence: one Crank--Nicolson step assembles a single tridiagonal
system $A\,T^{\,n+1}=d$, with two special rows carrying the physics
of \S\ref{sec:bc}. The surface row is a Dirichlet ghost (a
fixed-temperature entry) holding the Newton-solved skin temperature;
the basal row is a Neumann ghost (a fixed-flux entry) carrying \Qb. The Thomas algorithm solves the system in $O(N)$
operations, so a full lunation of half-hour steps costs only a linear
pass per step. Crank--Nicolson is unconditionally stable, so the
$\Delta t=1800$~s step is set by accuracy and certified by the ladder
of Appendix~\ref{app:numerics}, where the full stencil algebra, the
corrector sweeps for the temperature-dependent properties, and the
stability argument are recorded.

\begin{figure}[htbp]
  \centering
  \includegraphics[width=\linewidth]{figures/fig_numerics_grid_matrix.pdf}
  \caption[The discretization: grid, stencil, matrix]{The
  discretization in three views. \textbf{(a)}~The real geometric grid:
  cell width against cell-center depth on logarithmic axes
  ($\Delta z_0=2$~mm, growth $8\%$, 69 cells to $5$~m), with the
  anchor depth (dotted) and the deep-sensor band (shaded) marked.
  \textbf{(b)}~The finite-volume stencil: temperatures at cell
  centers, fluxes crossing faces at the harmonic-mean conductivity.
  \textbf{(c)}~The tridiagonal system one Crank--Nicolson step
  assembles: interior three-point rows, a Dirichlet ghost row at the
  surface (Newton-solved $T_s$), and a Neumann ghost row at the base
  carrying \Qb, solved by the Thomas algorithm in $O(N)$.}
  \label{fig:numerics}
\end{figure}

\section{Neglected forcing terms and their bounds}\label{sec:neglected}
The synodic cosine forcing omits the solar declination
($\pm1.54^{\circ}$) and the orbital eccentricity ($\pm3.3\%$
peak-to-peak in instantaneous flux). Because the retrieval compares
stability-window \emph{mean} temperatures rather than time-resolved
phases, only the cycle-mean insolation matters, and there the
neglected terms perturb at or below the percent level ($<0.1\%$ for
eccentricity). That is comparable to, and subsumed by, the albedo row
of the error budget (Table~\ref{tab:budget}).

The annual wave deserves its own sentence, because the deep sensors
sit within a meter of the surface. Its skin depth is
${\sim}30$--$40$~cm against $\lesssim10$~cm for the diurnal wave, so
at the retained depths ($\ge80$~cm, two or more annual skin depths
down) both waves have damped to negligible amplitude and the borestem
cut doubles as a dynamic-quiet cut. What the annual cycle does leave
behind is epochal: sensors whose stability windows sample different
calendar seasons carry slightly different means. That residual is
measured rather than modeled, by rebuilding every sensor's
$T_{\rm eq}$ on one shared 1974 window, and the result is carried as
the epoch-mixing row of the error budget
(\S\ref{sec:commonepoch}).

\section{Two-timescale structure of the problem}
\label{sec:twotimescale}
Figure~\ref{fig:skinwave} displays the structure that the whole
method exploits. Over one lunation the surface swings by
${\sim}280$~K, yet the swing collapses within the first half-meter.
Above that depth the column is a driven oscillator. Below it, the
column is effectively steady, carrying only the geothermal flux and a
rectified imprint of the oscillation. The diurnal skin and the
meter-scale sensors live on opposite sides of this divide, which is
why the deep sensors constrain \Kd\ cleanly, and why time-stepping
the whole $5$~m column is wasteful.

\begin{figure}[t]
  \centering
  \includegraphics[width=0.88\linewidth]{figures/fig_book_skinwave.pdf}
  \caption[The diurnal wave collapsing with depth]{One lunation of
  modeled temperature at increasing depths: the ${\sim}280$~K surface
  swing collapses within the first half-meter. Below the skin the
  column is quasi-steady.}
  \label{fig:skinwave}
\end{figure}

\section{Computational cost of brute-force integration}\label{sec:brute}
The observable is the \emph{periodic steady state}: the column that
repeats cycle after cycle once the drilling transient has decayed.
Reached naively, by time-stepping Eq.~\eqref{eq:heat} until the deep
column stops drifting, one solve costs about $3000$ simulated
lunations, because the deep relaxation time scales as $z^2/\kappa$.
The retrieval program of Chapter~\ref{ch:retrieval} needs tens of
thousands of solves, so the bootstrap alone would be computationally
prohibitive by brute force. The solver developed next reduces the cost
of one solve, and of the full sweep, by the factors quantified in
\S\ref{sec:cost}.

\section{The flux-anchored equilibrium method}\label{sec:overview}
Only the top few skin depths ever feel the diurnal forcing; below the
rectification zone the cycle-averaged column is already steady and
fully determined by the basal flux. The method exploits this split.
It time-steps \emph{only the skin} (Step~A), anchors the solution at
a depth $z_0=0.55$~m where the diurnal amplitude has fallen to
${\sim}0.15$~K, and reconstructs everything below the anchor by
integrating a one-variable steady-state ODE (Step~B). The two steps
alternate until the anchor temperature converges. The fixed point of
this composition is the periodic steady state of the full problem.

Figure~\ref{fig:methodflow} lays the loop out as a single cycle:
Step~A time-steps the skin and reads the cycle mean at the anchor,
Step~B integrates the closure ODE downward from there, and a drift
test either accepts the profile or feeds it back as the next guess.
In production the loop closes in about four passes.

\begin{figure}[t]
  \centering
  \includegraphics[width=\linewidth]{figures/fig_method_fluxanchored.pdf}
  \caption[The flux-anchored outer loop]{The flux-anchored outer loop.
  Step~A time-steps only the skin ($n_{\rm inner}$ lunations of
  Crank--Nicolson with a Newton surface balance and a Thomas solve),
  producing the cycle-mean skin, the anchor value $\Tm(z_0)$, and the
  rectified flux \urect. Step~B consumes those and reconstructs the
  deep column by integrating $d\Tm/dz=[\Qb-\urect]/K$ downward from
  the anchor. The reconstructed profile reseeds Step~A, and the loop
  stops when the anchor moves less than $5$~mK between cycles:
  typically four outer cycles (two on the coarse primer stage at
  $z_0=0.25$~m, two on the production stage at $z_0=0.55$~m).}
  \label{fig:methodflow}
\end{figure}

\section{Choice of the anchor depth}\label{sec:anchorpoint}
The method rests on one physical fact: there is a depth at which the
column has stopped oscillating but does not yet depend on the slow
interior relaxation. The diurnal wave, ${\sim}280$~K peak-to-peak at
the surface, is attenuated by $e^{-z/\delta}$ per skin depth $\delta$
of a few centimeters; by $z_0=0.55$~m its amplitude is
${\sim}0.15$~K, and the instantaneous profile at any phase of the
lunar day coincides with the cycle mean to within ${\sim}0.1$~K.
Above that depth the column is dynamic and must be time-stepped.
Below it the column is steady and can be computed. The anchor is the
hand-off between the two regimes: the deepest cell Step~A still
resolves honestly, and the initial condition from which Step~B
integrates downward.

The placement trades two costs and has a wide, certified optimum. The
anchor must sit below the rectification zone, where the covariance
correction \urect\ has decayed to under $1\%$ of \Qb, or the closure
of \S\ref{sec:closure} would depend on a quantity Step~A resolves
only approximately. It should sit no deeper than necessary, because
Step~A time-steps the truncated grid $[0,\,z_0+0.15~\mathrm{m}]$ and
every extra centimeter is paid on every inner lunation. Across
$z_0=0.45$--$0.70$~m the converged \Kds\ spans only $0.06$~\mWmK, a
half-range folded into the solver row of the error budget; the
reported values do not depend on the choice within that window.

In practice the solver approaches the anchor in two stages, the
structure visible in Fig.~\ref{fig:anchorconv}: a coarse stage
anchored at $0.25$~m (four inner lunations per cycle, $0.10$~K
tolerance) walks an arbitrary starting profile into the basin of
attraction cheaply, and the production stage at $z_0=0.55$~m
($n_{\rm inner}=96$, $5$~mK tolerance) finishes the job.

\section{Step B: the mean-flux closure ODE}\label{sec:closure}
Averaging Eq.~\eqref{eq:heat} over one lunation $P$ in periodic
steady state kills the storage term, because the column's enthalpy
returns to its starting value each cycle. The cycle-mean flux is then
constant in depth and pinned by the bottom boundary to \Qb. Writing
that mean flux in terms of the mean temperature introduces one
correction, the rectified flux
$\urect \equiv \langle K\,\partial_z T\rangle -
K(\Tm)\,\partial_z\Tm$, the covariance of the oscillating
conductivity and gradient, and yields the closure
\begin{equation}
  \boxed{\;\frac{d\Tm}{dz} = \frac{\Qb - \urect(z)}{K(\Tm, z)}\;}
  \label{eq:closure}
\end{equation}
(the sign convention is $\Qb-\urect$; Appendix~\ref{app:derivation}
gives every algebraic step). Above the anchor, \urect\ is evaluated
from the Step~A skin solution. Below it, \urect\ is negligible and
Eq.~\eqref{eq:closure} integrates downward as an ordinary
initial-value problem from the anchored temperature, by second-order
Runge--Kutta on the same geometric grid. Two variables $(z,t)$ have
become one $(z)$. That reduction is the entire speed-up.

The rectified flux deserves its own picture, because it is the least
intuitive object in the method. Figure~\ref{fig:urect} unpacks it at
Apollo~17. Panel~(a) shows the two oscillating factors at $15$~cm
depth over one lunation: the conductivity swings with temperature
through its radiative term, the gradient swings with the passing
thermal wave, and the two are correlated. Panel~(b) shows the
consequence: the mean of their product (the true cycle-mean flux,
which stays near \Qb) differs from the product of their means by
roughly seven times \Qb\ at this depth, and that gap is \urect. The
covariance of the oscillations carries heat that the mean-field
picture cannot see, which is exactly why the closure needs the
correction term. Panel~(c) shows why the anchor placement makes the
problem easy anyway: $|\urect|/\Qb$ collapses with depth and is
inside $1\%$ at the anchor, so below $z_0$ the closure reduces to
$d\Tm/dz\simeq\Qb/K$.

\begin{figure}[htbp]
  \centering
  \includegraphics[width=\linewidth]{figures/fig_urect_explainer.pdf}
  \caption[What the rectified flux is]{The rectified flux, unpacked
  from a converged Apollo~17 solve at the retrieved \Kds.
  \textbf{(a)}~At $15$~cm depth the conductivity (teal, left axis)
  and the temperature gradient (coral, right axis) oscillate together
  over the lunation. \textbf{(b)}~Because the factors are correlated,
  the mean of the product differs from the product of the means; the
  gap is \urect. The dashed line is \Qb. \textbf{(c)}~The magnitude
  of \urect\ relative to \Qb\ collapses with depth and is inside
  $1\%$ (dashed) at the anchor.}
  \label{fig:urect}
\end{figure}

One step of the walk makes the arithmetic concrete. At the Apollo~15
anchor the settled cycle-mean temperature is ${\sim}251.9$~K. There
the contact conductivity has fully reached its asymptote
($e^{-z_0/H}\approx10^{-4}$ at $H=6$~cm) and the radiative factor is
$1+\chi(\Tm/T_{\rm ref})^3\approx2.0$, so the total conductivity is
$K\approx2.0\times4.60\approx9.2$~\mWmK. The slope demanded by energy
conservation is $\Qb/K = 21/9.2 \approx 2.3$~K~m$^{-1}$, and the
first grid step below the anchor (${\sim}4$~cm there) adds
${\sim}0.10$~K. Repeating this node by node rebuilds the entire deep
column in a single downward pass, with no time-stepping and no
diffusion to wait for.

Figure~\ref{fig:anatomy} shows the anatomy of one converged solve.
Panel~(a) displays the mean column with its two constructions
colored: the time-stepped skin above the anchor and the reconstructed
deep column below it, joining seamlessly at the anchor point.
Panel~(b) quantifies the invariant the method is named after, as the
percent deviation of the cycle-mean flux from \Qb: large in the
rectification zone (off scale near the surface, where the rectified
component dominates), the deviation collapses through the skin,
enters the $\pm1\%$ band above the anchor, and stays within
$0.05\%$ of \Qb\ at and below it. That residual band is the
quantitative content of the claim that $\urect<1\%$ of \Qb\ below
the anchor.

\begin{figure}[htbp]
  \centering
  \includegraphics[width=\linewidth]{figures/fig_anchor_anatomy.pdf}
  \caption[Anatomy of one flux-anchored solve]{Anatomy of one
  converged solve (Apollo~17, retrieved \Kds). \textbf{(a)}~The mean
  temperature column, split into its two constructions: time-stepped
  skin (coral, Step~A) and reconstructed deep column (teal, Step~B),
  meeting at the anchor. \textbf{(b)}~The invariant, quantified: the
  percent deviation of the cycle-mean conductive flux from \Qb. The
  deviation is large in the rectification zone (off scale toward the
  surface), collapses through the skin, and stays inside the shaded
  $\pm1\%$ band from the anchor down.}
  \label{fig:anatomy}
\end{figure}

\section{Step A: skin-only time integration}\label{sec:stepA}
Step~A advances the full time-dependent equation, but only on a
truncated sub-grid reaching just below the anchor
($[0,\,z_0+0.15~\mathrm{m}]$); everything deeper is rebuilt by Step~B
each outer cycle. The truncation is structural, not an optimization.
When Step~A time-steps the full $5$~m column, the slow deep
relaxation leaks into the anchor and the retrieval acquires a
spurious dependence on the inner-loop length, the $n_{\rm inner}$
bias dissected in Appendix~\ref{app:numerics}. With the truncation,
the retrieved \Kds\ is independent of $n_{\rm inner}$ at the
certified settings and the convergence cost drops by a further factor
of about four.

\section{The outer loop as a fixed-point iteration}
\label{sec:fixedpoint}
The composition of Step~A on the skin with Step~B below the anchor is
a contraction. Step~A relaxes any starting skin toward the true skin
solution within a few lunations, because the skin's diffusive
timescale is short; Step~B is exact below the anchor given a correct
anchor temperature. The loop therefore converges geometrically: in
production the anchor drift collapses from $46$~mK on the
shallow-anchor stage to $0.4$~mK at the $0.55$~m anchor, entering the
$\pm5$~mK tolerance at the fourth outer cycle
(Fig.~\ref{fig:anchorconv}). Convergence is judged on the anchor
drift, never on a fixed cycle count; the fixed-count trap and its
consequences are part of the certification record
(Appendix~\ref{app:numerics}).

\begin{figure}[htbp]
  \centering
  \includegraphics[width=0.95\linewidth]{figures/fig_anchor_convergence.pdf}
  \caption[Outer-loop convergence of the anchor temperature]{The
  two-stage anchor settling within four outer cycles (top) and the
  cycle-to-cycle drift falling below both stage tolerances (bottom),
  from the real production trace.}
  \label{fig:anchorconv}
\end{figure}

The construction, its brute-force verification, and the invariant it
rests on are collected in Fig.~\ref{fig:anchorfilm}. Panel~(a) builds
the deep column outward from the anchor. Panel~(b) lays that
reconstruction over a full brute-force solve; the two curves fall on
top of each other. Panel~(c) shows the invariant itself, the
cycle-mean flux sitting flat at \Qb\ everywhere below the skin.

\begin{figure}[htbp]
  \centering
  \includegraphics[width=\linewidth]{figures/fig_anchor_method.pdf}
  \caption[The flux-anchored method: construction, correctness,
  invariant]{The flux-anchored method at the retrieved Apollo~15
  $\Kds=4.60$~\mWmK. \textbf{(a)}~Construction: from a deliberately
  wrong initial profile, the shallow column is time-stepped until its
  cycle mean settles (green); the deep column (teal) is reconstructed
  by walking down from the anchor at $z_0=0.55$~m with the slope rule
  $d\Tm/dz=[\Qb-\urect]/K$, landing on the converged steady state
  (dashed, coincident). \textbf{(b)}~Correctness: a brute-force
  spin-up started $13$~K wrong approaches the anchored solution over
  ${\sim}10^{3}$ lunations, closing to $33$~mK after $3000$ cycles;
  the anchored solve reaches the same state in four outer cycles.
  \textbf{(c)}~The invariant the reconstruction integrates: the total
  cycle-mean conductive flux equals \Qb\ at every depth below the
  rectification zone.}
  \label{fig:anchorfilm}
\end{figure}

\section{Measured computational cost}\label{sec:cost}
Cost is reported as a hardware-independent \emph{speed-up factor}
rather than as wall-clock time, since the absolute time per solve
depends on the machine while the ratio between two methods run on the
same machine does not. Table~\ref{tab:speed} and Fig.~\ref{fig:speedup}
give those factors (each the median of three timed runs); the absolute
timings behind them, and the machine on which they were measured, are
recorded in Appendix~\ref{app:timings} for concreteness. The
algorithmic gain, anchored against brute force with both interpreted,
is about $21\times$ per solve. Compiling the Step-A march (a kernel
checked against a verified C++ reference) multiplies this by a further
${\sim}117\times$, for a total of about $2500\times$ over the
interpreted brute-force solve, and the same factor carries to the full
61-point sweep. The remaining gap between the compiled Python kernel
and a line-for-line C++ reference is a constant factor of about
$1.4\times$ per lunation, not an algorithmic difference, and all three
implementations agree to ${\sim}1.1\times10^{-12}$~K.

\begin{table}[htbp]
  \centering\small
  \caption{Speed-up of the flux-anchored solver, expressed as
  hardware-independent factors relative to the interpreted brute-force
  solve (each the median of three timed runs).}
  \label{tab:speed}
  \begin{tabular}{lc}
    \toprule
    Stage & Speed-up factor \\
    \midrule
    Flux-anchored vs.\ brute force (both interpreted) & ${\sim}21\times$ \\
    Compiled kernel vs.\ interpreted (Step~A)         & ${\sim}117\times$ \\
    Total (compiled anchored vs.\ interpreted brute force) & ${\sim}2500\times$ \\
    \midrule
    C++ reference vs.\ compiled Python kernel (per lunation) & ${\sim}1.4\times$ \\
    \bottomrule
  \end{tabular}
\end{table}

\begin{figure}[t]
  \centering
  \includegraphics[width=0.82\linewidth]{figures/fig_speedup_factors.pdf}
  \caption[Solver speed-up]{Speed-up of the flux-anchored solver
  relative to the interpreted brute-force solve, per solve and for the
  full 61-point sweep, on a logarithmic factor axis. Each bar is the
  median of three timed runs; the factors are ratios and do not depend
  on the machine.}
  \label{fig:speedup}
\end{figure}

\section{Certification summary}\label{sec:certsummary}
Speed bought with a wrong answer is worthless, so equality with the
ground truth is certified rather than assumed, at three levels.
First, a fully converged brute-force run reaches the same profile
(Fig.~\ref{fig:anchorfilm}b). Second, the discretization was pushed
down a time-step ladder until the \emph{retrieved} \Kds, the quantity
actually reported, stopped moving; that process exposed and
eliminated a first-order defect in the time march (the missing wrap
step). Third, an independent C++ implementation reproduces the solver
to $10^{-12}$~K. The full record, including the three input errors
found along the way and their effect on the headline number, is
Appendix~\ref{app:numerics}.

\section*{Chapter summary}
The physics of this chapter reduces to one split and one invariant.
The split: a ${\sim}280$~K diurnal oscillation confined to the top
half-meter sits above a deep column that is steady on the lunation
timescale, so only the skin ever needs time-stepping. The invariant:
in periodic steady state the cycle-mean conductive flux equals \Qb\
at every depth, so the deep mean profile follows from a one-variable
ODE anchored at $z_0=0.55$~m. Exploiting the pair cuts the cost of one
solve by about a factor of $2500$ at certified-equal answers, and that
speed-up is what makes the next chapter's uncertainty program (a
1500-draw bootstrap, a mapped degeneracy, MCMC cross-checks)
affordable at all. Every
knob the solver introduces (anchor depth, inner length, time step,
grid) has a demonstrated convergence plateau in the retrieved \Kds\
itself, recorded in Appendix~\ref{app:numerics}.
Table~\ref{tab:config} collects the certified production
configuration in one place for reference.

\begin{table}[h!]
  \centering\small
  \caption[The certified production configuration]{The certified
  production configuration of the solver and retrieval. Every entry
  has a demonstrated \Kds\ convergence plateau or sensitivity sweep
  behind it (Appendix~\ref{app:numerics}; \S\ref{sec:budget}).}
  \label{tab:config}
  \begin{tabular}{lll}
    \toprule
    Setting & Value & Certification \\
    \midrule
    Spatial grid & $\Delta z_0=2$~mm, growth $8\%$, 69 cells to $5$~m & grid-halving re-retrieval \\
    Time step & $\Delta t=1800$~s, 2 corrector sweeps & $\Delta t$ ladder on \Kds \\
    Anchor & $z_0=0.55$~m, tolerance $5$~mK & $z_0=0.45$--$0.70$~m sweep \\
    Primer stage & $z_0=0.25$~m, tolerance $0.10$~K & basin-of-attraction check \\
    Inner length & $n_{\rm inner}=96$ lunations & plateau at $64/96/128$ \\
    Step-A sub-grid & $[0,\,z_0+0.15~\mathrm{m}]$ & vs.\ full column, $2\times10^{-3}$ \\
    \Kd\ sweep & 29 points (A15), 32 points (A17) & vertex vs.\ denser grids \\
    Bootstrap & 1500 resamples, $\pm2.5$~cm jitter & BCa comparison \\
    \bottomrule
  \end{tabular}
\end{table}

% ============================================================================
\chapter{Retrieval and uncertainty}\label{ch:retrieval}

\section{Objective function and parameter sweep}\label{sec:sweep-method}
The retrieval minimizes the deep-sensor misfit
\begin{equation}
  \mathrm{RMSE}(\Kd) = \Bigl\langle\bigl[T_{\rm model}(z_i;\Kd)
    - T_{\rm eq}(z_i)\bigr]^2\Bigr\rangle_{z_i\ge 80~\mathrm{cm}}^{1/2},
  \label{eq:rmse}
\end{equation}
with every other model input held at its published value
(Table~\ref{tab:params}); the RMSE (root-mean-square error) is the
typical size, in kelvin, of the model-minus-data mismatch over the
retained sensors. For each site, \Kd\ is swept across a dense
grid (29 points spanning $1.0$--$15.0$~\mWmK\ at Apollo~15, 32 points
spanning $3.0$--$25.0$~\mWmK\ at Apollo~17, refined near the minimum),
each grid point solved to periodic equilibrium, and the model column
interpolated to the sensor depths. The retrieved \Kds\ is the vertex
of a parabola through the three points bracketing the minimum. Grid
convergence of the vertex was verified against coarser and denser
grids; a uniform coarse grid biases the vertex by ${\sim}0.04$~\mWmK,
one of the certification items of Appendix~\ref{app:numerics}. The
retrieval has one free parameter per site.

\section{Statistical uncertainty: the bootstrap}\label{sec:bootstrap}
The primary uncertainty estimate is a non-parametric bootstrap
\citep{efron1993,press2007}. The idea is simple: rebuild the sensor
set many times with the random variations the data themselves allow,
refit each version in full, and read the spread of the refitted
answers as the statistical uncertainty. Concretely, 1500 resamples of
the deep-sensor set are drawn with replacement, each with sensor
depths jittered by the $\pm2.5$~cm placement uncertainty
\citep{nagihara2018}, and each is re-retrieved through the full
sweep. The bootstrap propagates exactly
the two quantities measured at the boreholes, which sensors exist and
where they sit, and assumes nothing about the basal flux. An analysis
of this size only became affordable with the fast solver.

Figure~\ref{fig:bootdraw} makes one of the 1500 draws concrete. In
panel~(a) the redraw doubles the $84$~cm sensor, omits the deepest,
and shifts every depth within its placement uncertainty; refitting
that set end to end contributes one draw to panel~(b)'s distribution.
The histogram's shape is itself informative: its long upper tail is
the flat high-\Kd\ shoulder of the Apollo~15 misfit bowl showing up
as probability mass, which is where the asymmetric
$^{+0.86}_{-0.21}$ interval comes from.

\begin{figure}[htbp]
  \centering
  \includegraphics[width=0.9\linewidth]{figures/fig_bootstrap_draw.pdf}
  \caption[One bootstrap draw, made concrete]{One bootstrap draw made
  concrete at Apollo~15. \textbf{(a)}~The original deep-sensor set
  (open) and one with-replacement redraw (coral): one sensor drawn
  twice, one not drawn, every depth jittered within the $\pm2.5$~cm
  placement uncertainty. \textbf{(b)}~The real distribution of 1500
  such redraws, each refit end to end, with the point estimate and
  the $95\%$ bounds marked.}
  \label{fig:bootdraw}
\end{figure}

Percentile intervals are reported throughout. Bias-corrected
accelerated (BCa) intervals agree with the percentile endpoints to
within $1.3$~\mWmK\ at both sites and on the contrast, a difference
comparable to the interval half-widths themselves at these sample
sizes ($N=7$ and $16$), so the intervals should be read as
approximate rather than exact coverage statements.

One structural limitation belongs here rather than in the discussion.
Each site's sensors ride on only two probes, and sensors sharing a
borestem share an emplacement disturbance and a calibration chain;
Table~\ref{tab:sensors} shows the Apollo~17 probe-2 string reading
systematically warmer than probe~1 at matched depths. The bootstrap
treats the $7$ and $16$ sensors as exchangeable, so the effective
number of independent units is smaller than the sensor count and the
quoted intervals are best read as lower bounds on the true
uncertainty. A probe-clustered (hierarchical) resampling scheme would
widen them; it is identified as follow-on work in
\S\ref{sec:caveats}.

\section{Model selection by information criterion}\label{sec:aicc-method}
Models with different parameter counts are compared on the same
sensor set with the small-sample-corrected Akaike information
criterion \citep{burnham2002}. A reduced $\chi^2$ is deliberately not
used: the within-window scatter ($0.03$--$0.05$~K) measures
short-term thermometer stability, not model adequacy, and would give
$\chi^2_\nu\gg1$ for every model, informative about none. Following
\citet{burnham2002}, $\Delta\mathrm{AICc}\lesssim2$ means
statistically indistinguishable models and
$\Delta\mathrm{AICc}\gtrsim10$ a decisive preference.

\section{Provenance and classification of inputs}\label{sec:register}
A retrieval is only as defensible as its inputs, so every input is
classified by the evidence behind it and treated accordingly
(Table~\ref{tab:register}). Published measurements carry their
literature envelope into the error budget. Numerical knobs must show
a demonstrated convergence plateau in \Kds. Discretionary choices are
swept and reported as sensitivity rows. Derived quantities carry
their derivation.

\begin{table}[htbp]
  \centering\small
  \caption{The input register: every input to the headline number, by
  evidence class.}
  \label{tab:register}
  \begin{tabular}{lll}
    \toprule
    Class & Inputs & Treatment \\
    \midrule
    Published measurement & \Qb, \Ks, $H$, $\chi$, $\rho(z)$, $c_p(T)$, & envelope $\to$ error \\
      & albedo, emissivity, $S_0$, sensor depths & budget row \\
    Numerical knob & $\Delta t$, $n_{\rm inner}$, spatial grid, \Kd-grid & convergence plateau \\
      & & certified (App.~\ref{app:numerics}) \\
    Discretionary choice & $80$~cm cut, anchor $z_0$, stability & swept; sensitivity \\
      & threshold & row in budget \\
    Derived & $z_0$ below the rectification zone & derivation shown \\
    \bottomrule
  \end{tabular}
\end{table}

\section{Systematic error budget}\label{sec:budget}
Every fixed input is perturbed through its literature envelope and
the retrieval repeated; Table~\ref{tab:budget} summarizes. Each row
answers one concrete question: if this input is wrong by its published
uncertainty, how far does the retrieved conductivity move. Although
the systematic rows dominate it, the quadrature total also folds in
the statistical row, so the table is the full error budget rather
than a systematics-only figure. At
Apollo~15 the dominant systematic is \Qb\ itself, evaluated on the
direct degeneracy map of \S\ref{sec:robustness} rather than on an
assumed scaling. At Apollo~17 the surface albedo overtakes it as the
single largest term in the budget, because part of the sweep range
there responds nearly flatly to \Kd\ and even a $\pm0.01$ albedo
change moves the minimum substantially; the basal flux and
surface-conductivity terms remain material. The quadrature totals are
$\sigma=1.88$ (A15) and $3.88$ (A17)~\mWmK. The A17 total includes
the epoch-mixing term, the largest single methodological systematic
there (\S\ref{sec:commonepoch}). Two rows are conditional rather than
summed: the radiative coefficient $\chi$ and the compaction scale $H$
parameterize the functional form the retrieval holds fixed, so \Kds\
is reported as conditional on their published values
(\S\ref{sec:caveats}).

The budget is reported at $1\sigma$, and a wider $3\sigma$ envelope
would be wrong here. Most of the systematic rows are half-ranges of
the retrieval across a bounded literature interval, not standard
deviations of anything Gaussian; multiplying them by three would
claim the input can sit three times outside its published range, and
it drives $\Kds(\mathrm{A15})$ below zero, an unphysical
conductivity. The size of the budget is fixed by how well the basal
flux and albedo are known, not by the coverage level chosen for
reporting. Where the scatter is genuinely Gaussian, in the bootstrap,
the statistical part alone clears $3\sigma$ against the global value
(A15 $[3.7,\,7.4]$, A17 $[5.8,\,8.5]$~\mWmK). Only the inter-site
contrast fails a $3\sigma$ test, which is why it is reported as
marginal throughout.

\begin{table}[htbp]
  \centering\small
  \caption{Error budget for \Kds\ (\mWmK). Values are
  half-ranges of the retrieval across each input's literature
  envelope; the \Qb\ row uses the direct degeneracy map of
  \S\ref{sec:robustness}.}
  \label{tab:budget}
  \begin{tabular}{lcc}
    \toprule
    Source & A15 & A17 \\
    \midrule
    Statistical (bootstrap $1\sigma$) & 0.53 & 0.52 \\
    Basal flux \Qb                    & 1.61 & 1.66 \\
    Albedo ($\pm0.01$)                & 0.72 & 2.96 \\
    Surface conductivity \Ks\ ($\pm30\%$) & 0.36 & 1.55 \\
    Solver/discretization             & 0.04 & 0.04 \\
    Sensor-depth / threshold / density (each) & $\le$0.12 & $\le$0.33 \\
    Epoch mixing (common-1974 re-fit) & $<$0.05 & 0.89 \\
    \midrule
    \textbf{Total (quadrature)}       & \textbf{1.88} & \textbf{3.88} \\
    \midrule
    Conditional: $H$ ($3$--$10$~cm)   & 0.51 & 3.47 \\
    Conditional: $\chi$               & \multicolumn{2}{c}{reported conditional on $\chi=2.7$} \\
    \bottomrule
  \end{tabular}
\end{table}

One honest qualifier accompanies the totals. The statistical and \Qb\
rows both propagate along the same $\Kd$--$\Qb$ ridge, so the
quadrature assumes an independence that is only approximate, and
several rows are half-ranges of visibly asymmetric responses (the A17
albedo row especially). The totals are sensitivity summaries,
reliable for judging which inputs matter, not calibrated standard
errors.

\subsection*{A note on the two uncertainties}
Because two different uncertainties travel with each retrieved value,
it is worth stating plainly how they relate, so they are never
conflated. The \emph{statistical} uncertainty is the random scatter
that comes from having only a handful of noisy sensors; it is
estimated by the bootstrap (\S\ref{sec:bootstrap}) and is small. The
\emph{systematic} uncertainty is the bias that would follow if a fixed
input, chiefly the basal heat flux, were wrong within its published
range; it does not shrink with more sensors and, as
Table~\ref{tab:budget} shows, it is the larger of the two. Both are
quoted at $1\sigma$, the standard convention for a measured value,
which spans roughly the central $68\%$ of the distribution; the wider
$95\%$ interval, about $\pm2\sigma$, is the one listed in brackets in
Eq.~\eqref{eq:headline} and the results tables. A $1\sigma$ range is
therefore narrower than a $95\%$ one not because the error has been
made to look small but because the two describe the same scatter at
different confidence levels. The single honest uncertainty on each
\Kds\ is the \emph{total}, $\pm1.9$ (A15) and $\pm3.9$ (A17)~\mWmK,
which the systematic dominates. Where a significance claim is made
rather than a value reported, in particular the inter-site contrast of
Chapter~\ref{ch:results}, the stricter $95\%$ level is used, and the
contrast is reported as marginal.

\section*{Chapter summary}
The retrieval machinery is deliberately narrow: one free parameter
per site, an RMSE objective over the retained sensors, and a
parabolic vertex on a converged sweep. Around that narrow core sits
the uncertainty apparatus that carries the thesis. The bootstrap
propagates what the boreholes actually measured (sensor-set
composition and placement) and gives the statistical intervals of
Eq.~\eqref{eq:headline}; the AICc handles model comparison where a
reduced $\chi^2$ would be uninformative; the input register
classifies every fixed quantity by its evidence; and the systematic
budget converts each published envelope into a \Kds\ half-range,
totalling $\pm1.9$ and $\pm3.9$~\mWmK\ with the basal flux and the
Apollo~17 albedo dominant. The next chapter reports what this
machinery finds.

% ============================================================================
\chapter{Results and robustness}\label{ch:results}

\section{Retrieved conductivities}\label{sec:headline-results}
The deep-sensor RMSE has a well-resolved minimum at each site:
\begin{equation}
  \Kds(\mathrm{A15}) = 4.60^{+0.86}_{-0.21}\,[4.18,\,6.96], \qquad
  \Kds(\mathrm{A17}) = 7.08^{+0.53}_{-0.52}\,[6.16,\,8.07]\ \mWmK
  \label{eq:headline}
\end{equation}
($1\sigma$ asymmetric about the point estimate; brackets $95\%$
bootstrap). Figure~\ref{fig:headline} shows the principal result at a
glance: the two per-site retrievals with their intervals, and the
global $3.4$~\mWmK\ lying outside both. Neither interval reaches down
to the global value, and the two barely overlap, so the figure
carries both claims at once. Each site is more conductive than the
Moon-wide $3.4$, and Apollo~17 is more conductive than Apollo~15. The
overlap that remains, at the top of the Apollo~15 interval and the
bottom of the Apollo~17 interval, is where the contrast loses its
statistical margin.

\begin{figure}[htbp]
  \centering
  \includegraphics[width=0.85\linewidth]{figures/fig_headline.pdf}
  \caption[Per-site \Kds\ with 95\% intervals]{Per-site \Kds\ with
  95\% bootstrap intervals against the \citet{hayne2017} global
  value.}
  \label{fig:headline}
\end{figure}

The sweep curves behind those numbers are shown in
Fig.~\ref{fig:sweep}. The two bowls differ in shape, and the shape is
the uncertainty. The Apollo~15 bowl is shallow and nearly flat toward
high \Kd, so the fit tolerates much larger values more readily than
smaller ones and its interval runs far upward but little downward.
The deeper, more symmetric Apollo~17 bowl pins its value tightly.
The two bowls also sit at noticeably different heights, and the height
is just the best-fit residual: the Apollo~15 curve floats higher
because even its best \Kd\ still leaves a $1.0$~K mismatch (only seven
sensors in a narrow band, with more probe-to-probe scatter), whereas
the Apollo~17 fit reaches $0.4$~K on sixteen more consistent sensors.
The bowl's depth measures how cleanly the site's data can be fit; its
shape measures how well \Kd\ is pinned.
Both minima sit well to the right of the Hayne global $3.4$ and the
\citet{feng2020} deep value of $3.8$, so the departure from the
global number can be read straight off the curves. Extending the
sweep past the fitting grid shows each curve rising again on the
high-\Kd\ side: the minima are genuine bracketed bowls, not
descending shoulders.

\begin{figure}[htbp]
  \centering
  \includegraphics[width=0.9\linewidth]{figures/fig_kd_sweep.pdf}
  \caption[Deep-sensor RMSE against \Kd\ at both sites]{Deep-sensor
  RMSE as a function of \Kd\ at both sites. Stars: retrieved \Kds;
  horizontal bars: 95\% bootstrap intervals; vertical lines: the
  \citet{hayne2017} global value and the \citet{feng2020} deep
  value.}
  \label{fig:sweep}
\end{figure}

Figure~\ref{fig:sweepworked} walks the Apollo~17 case as a worked
example: the bowl from the real sweep, the three grid points that
bracket its minimum, and the parabola vertex that defines
$\Kds=7.08$~\mWmK.

\begin{figure}[htbp]
  \centering
  \includegraphics[width=0.85\linewidth]{figures/fig_sweep_worked.pdf}
  \caption[The A17 retrieval as a worked example]{The certified
  Apollo~17 retrieval as a worked example: the RMSE bowl, the three
  bracketing grid points, and the parabolic vertex at
  $\Kds=7.08$~\mWmK.}
  \label{fig:sweepworked}
\end{figure}

The bootstrap distributions are shown in Fig.~\ref{fig:boot}.
Panel~(a) shows why the point estimates look clean: each site's
redraws pile into a narrow peak, and the two peaks barely touch. The
Apollo~15 peak carries a long tail to the right, inherited directly
from the flat high-\Kd\ side of its bowl in Fig.~\ref{fig:sweep}: many
resampled sensor sets place the minimum at a high \Kd, since the fit
barely worsens there, while the steep low-\Kd\ side keeps the left edge
sharp, which is exactly the asymmetry ($^{+0.86}_{-0.21}$) reported for
the value. The rounder Apollo~17 bowl gives a nearly symmetric peak.
Panel~(b) is the honest counterweight. The inter-site contrast,
formed by pairing bootstrap draws, has median $+2.31$~\mWmK, $95\%$
CI $[-0.12,\,3.56]$, and one-sided tail
$P_{\rm boot}(\Delta\Kd\le0)=0.031$. The distribution sits mostly
above zero but leaves a thin tail across it; those crossing draws,
the $3.1\%$ that place Apollo~15 above Apollo~17, are why the
magnitude is only marginally resolved even though the sign rarely
flips. (The median of paired differences, $2.31$, differs slightly
from the difference of the point estimates, $2.48$, by
construction.)

\begin{figure}[htbp]
  \centering
  \includegraphics[width=0.95\linewidth]{figures/fig_bootstrap.pdf}
  \caption[Bootstrap distributions and the contrast]{\textbf{(a)}~Per-site
  bootstrap distributions of \Kds\ (1500 draws each) with the
  \citet{hayne2017} global value marked; \textbf{(b)}~the inter-site
  contrast distribution with its 95\% interval and one-sided tail
  $P_{\rm boot}=0.031$.}
  \label{fig:boot}
\end{figure}

The fitted columns against the data are shown in
Fig.~\ref{fig:profiles}. The global-\Kd\ forward model misses the
deep sensors by $1.09$~K at Apollo~15 and $0.89$~K at Apollo~17, with
warm biases of $+0.62$ and $+0.54$~K; the site fits reduce the RMSE
to $1.00$ and $0.40$~K and the biases to $+0.44$ and $+0.12$~K. The
Apollo~17 improvement is dramatic and visually obvious in the
deep-sensor zoom. The Apollo~15 improvement is modest in RMSE but
moves the profile's \emph{slope} into agreement, and the deep
gradient is the physically loaded feature.

\begin{figure}[htbp]
  \centering
  \includegraphics[width=\linewidth]{figures/fig_thermal_profiles.pdf}
  \caption[Modeled columns vs.\ the deep sensors]{Modeled temperature
  columns against the HFE sensors at both sites (top), with
  deep-sensor zooms (bottom). Site-fitted Hayne profiles (solid)
  against the \citet{martinez2021} forward evaluation (dashed).}
  \label{fig:profiles}
\end{figure}

\section{Model selection}\label{sec:modelsel}
Model selection asks whether the extra freedom of a per-site \Kd\
earns its keep or merely fits noise. The tool is the
small-sample-corrected Akaike information criterion (AICc), which
scores a model by adding a penalty for each free parameter to its
misfit: a lower AICc rewards a better fit but charges for complexity,
so it favors the simpler model unless the fit improvement is large
enough to justify the added parameter. Only differences between models
matter, not the absolute value; a gap $\Delta\mathrm{AICc}\lesssim2$
means two models are effectively tied, while
$\Delta\mathrm{AICc}\gtrsim10$ is a decisive preference.

The AICc accounting supports the site-resolved description honestly
rather than uniformly (Table~\ref{tab:aicc}). At Apollo~17 the site
fit is decisively preferred over the global value
($\Delta\mathrm{AICc}\approx-23$). At Apollo~15, where the RMSE
improvement is $8\%$, parsimony alone would keep the global value
($\Delta\mathrm{AICc}\approx+3$ once the parameter penalty is paid);
the Apollo~15 departure rests on the bootstrap interval, whose lower
bound ($4.18$) sits above $3.4$. Jointly the two-value description is
preferred, with $\Delta\mathrm{AICc}\approx-20$, carried almost entirely
by Apollo~17. A discrete three-layer alternative fitted per site does
not outperform the site-fitted Hayne form at either site.

\begin{table}[htbp]
  \centering\small
  \caption{Model selection on the deep sensors. Lower AICc is better;
  differences within a site are the meaningful quantity.}
  \label{tab:aicc}
  \begin{tabular}{lcccc}
    \toprule
    & \multicolumn{2}{c}{A15} & \multicolumn{2}{c}{A17} \\
    Model & RMSE (K) & AICc & RMSE (K) & AICc \\
    \midrule
    Hayne, global $\Kd=3.4$ & 1.09 & 4.1 & 0.89 & $-1.4$ \\
    Hayne, site-fit \Kds    & 1.00 & 7.0 & 0.40 & $-24.6$ \\
    3-layer, site-fit       & 0.98 & 6.7 & 0.92 & 2.4 \\
    \bottomrule
  \end{tabular}
\end{table}

Figure~\ref{fig:aiccanat} shows why the criterion splits the two
sites, by separating its two ingredients. At Apollo~15 the $8\%$ RMSE
gain is worth only $-1.2$ of fit improvement, while adding a
parameter at $N=7$ costs $+4.2$ of small-sample penalty; the net
$+3.0$ is parsimony's verdict. At Apollo~17 the halved RMSE is worth
$-25.6$ against a $+2.6$ penalty, and the net $-23$ sits far past the
decisive threshold. Same criterion, same honesty, opposite verdicts,
driven by how much residual each site's refit actually removes.

\begin{figure}[htbp]
  \centering
  \includegraphics[width=0.66\linewidth]{figures/fig_aicc_anatomy.pdf}
  \caption[AICc decomposed at the two sites]{The AICc comparison of
  Table~\ref{tab:aicc} decomposed into its two ingredients: the fit
  improvement $N\ln(\mathrm{RSS_{fit}}/\mathrm{RSS_{glob}})$ (teal)
  against the small-sample parameter penalty (coral), with the net
  (diamonds) and the $\Delta\mathrm{AICc}=-10$ decisive threshold
  dashed. Component values follow from the RMSEs of
  Table~\ref{tab:aicc}; totals match its differences within
  rounding.}
  \label{fig:aiccanat}
\end{figure}

\section{Internal validation}\label{sec:holdout}
Two held-out checks confirm the retrieval is not driven by any single
sensor or sensor type (Fig.~\ref{fig:holdout}). Fitting on
gradient-bridge sensors alone and predicting the ring-bridge sensors,
and vice versa, reproduces held-out temperatures within
${\sim}0.11$~K of the in-sample values at both sites. Withholding the
deepest sensor at each site ($139$~cm at Apollo~15, $234$~cm at
Apollo~17) and refitting on the remainder predicts the held-out
temperature to $0.37$ and $0.25$~K. Both are internal checks: they
establish that the fit generalizes across thermometer families and
depths, not that the inter-site contrast is externally confirmed.

\begin{figure}[htbp]
  \centering
  \includegraphics[width=\linewidth]{figures/fig_holdout.pdf}
  \caption[Held-out validation]{Held-out validation.
  \textbf{(a)}~TG/TR cross-prediction RMSEs against in-sample values;
  \textbf{(b)}~leave-one-deepest-out prediction residuals.}
  \label{fig:holdout}
\end{figure}

\section{Common-epoch sensitivity}\label{sec:commonepoch}
The stability windows end at different calendar epochs while the
record still warms, and the mixing is not random. At Apollo~17 the
four retained TG sensors at $130$--$178$~cm carry mid-1977 windows
while the interleaved TR sensors and every sensor below $180$~cm
carry 1974 windows, a $3.1$-year spread aligned with depth; at
Apollo~15 the spread is $1.4$ years with trailing slopes at or below
$0.22$~K~yr$^{-1}$. Figure~\ref{fig:epochmap} makes the mixing visible: each retained
sensor's window drawn on a shared calendar axis, ordered by depth.
The two shallow TG pairs sit isolated on the right, their windows
running to mid-1977, while every other bar ends with 1974, and the
split aligns exactly with depth. Windows this cleanly separated in
time, on a record still warming, are what make the epoch systematic
worth measuring rather than assuming away.

\begin{figure}[htbp]
  \centering
  \includegraphics[width=0.88\linewidth]{figures/fig_epoch_map.pdf}
  \caption[The Apollo 17 window epochs by depth]{The Apollo~17
  stability windows on a shared day axis, ordered by sensor depth
  (charcoal TG, gray TR; co-depth pairs overplot). The shallow TG
  pairs carry mid-1977 windows while the TR pairs and every sensor
  below $180$~cm end with 1974, a $3.1$-year spread aligned with
  depth. This depth-aligned epoch mixing is the systematic measured
  by the common-epoch reduction.}
  \label{fig:epochmap}
\end{figure}

An epoch-aligned reduction tests what this
mixing costs. Rebuilding every retained sensor's $T_{\rm eq}$ from
one shared 1974 calendar window (every sensor has data there, so no
extrapolation is involved) leaves Apollo~15 unchanged
($|\Delta\Kds|\le0.02$) but moves Apollo~17 to $\Kds=6.16$~\mWmK\
(bootstrap $95\%$ CI $[5.92,\,6.81]$), a shift of $-0.89$, and
improves the fit (RMSE $0.40\to0.36$~K), as expected if the
mixed-epoch profile was internally inconsistent. A slope-translated
construction agrees to $0.03$~\mWmK.

Under the common epoch the contrast is $+1.55$~\mWmK\ ($95\%$ CI
$[-0.11,\,2.31]$, $P_{\rm boot}=0.04$): still positive, still
rejecting the global value at both sites, and more marginal than at
the certified windows. Epoch mixing is therefore the largest single
methodological systematic at Apollo~17 and is carried as the
epoch-mixing row of Table~\ref{tab:budget}. Both A17 values, $7.08$
at the certified windows and $6.16$ at the common epoch, are quoted
through the discussion.

The certified windows remain the primary reduction, and the choice
deserves an explicit defense, because it is not the conservative one
for the contrast claim. The stability-window rule is the pre-stated
estimator of this thesis, applied uniformly to every sensor: each
certified window is the longest tail consistent with equilibrium for
that sensor, which minimizes its individual drift bias. The shared
1974 window buys cross-sensor epoch consistency by truncating the
four shallow Apollo~17 TG records to an earlier, still-relaxing
portion of their histories, so their $T_{\rm eq}$ estimates carry
more residual emplacement drift, not less; its better RMSE reflects
internal consistency of the profile, not better individual
temperatures. Choosing the common epoch instead would lower both the
Apollo~17 value and the contrast, which is exactly why both
reductions are reported side by side and why the contrast is called
marginal under either.

\section{Comparison with an independent global model}\label{sec:secondmodel}
An independent forward evaluation of the \citet{martinez2021}
$K(T,\rho)$ conductivity model, at its published coefficients with no
per-site freedom, misses the deep sensors by $1.05$~K at Apollo~15
and $0.63$~K at Apollo~17, warm in the same direction as the global
Hayne value. Because that model carries an explicit density
dependence, it admits a sharper diagnostic than a scalar refit: sweep
a per-site density scalar $\alpha$, with
$\rho_d^{\rm site}=\alpha\cdot1800$~kg\,m$^{-3}$ entering the
Mart\'inez coefficients while everything else stays fixed, and ask
where each site's misfit bottoms out (Fig.~\ref{fig:alphasweep}).

Both RMSE($\alpha$) bowls bottom inside or at the edge of the
Apollo-core admissible density band ($1700$--$2000$~kg\,m$^{-3}$;
\citealp{mitchell1973,carrier1991}): $\alpha^*_{\rm A15}=1.19$
($\rho_d^*\approx2139$~kg\,m$^{-3}$, RMSE $0.91$~K) and
$\alpha^*_{\rm A17}=1.04$ ($\rho_d^*\approx1877$~kg\,m$^{-3}$, RMSE
$0.60$~K). The published $K(T,\rho)$ is therefore physically
admissible at both boreholes. The flat Apollo~17 bowl also exposes
the limit: no admissible density closes the remaining gap to the
per-site Hayne fit ($0.60$ against $0.40$~K), so sharpening the A17
gradient further requires freedom in the conductivity mechanism, the
contact/radiative split, and not in the density input. Two
independently calibrated global forms miss the same sensors in the
same direction, which locates the deficiency in globally averaged
calibration rather than in either functional form.

\begin{figure}[htbp]
  \centering
  \includegraphics[width=0.9\linewidth]{figures/fig_alpha_sweep.pdf}
  \caption[Per-site density diagnostic under the Mart\'inez
  form]{Diagnostic per-site retrieval under the \citet{martinez2021}
  $K(T,\rho)$ form: deep-sensor RMSE versus the per-site density
  scalar $\alpha$ ($\rho_d^{\rm site}=\alpha\cdot1800$~kg\,m$^{-3}$),
  with the retrieved $\alpha^*$ starred, the published baseline at
  $\alpha=1$ dashed, the Apollo-core admissible band
  ($1700$--$2000$~kg\,m$^{-3}$) marked by verticals, and the
  solid-basalt bound ($3000$~kg\,m$^{-3}$; \citealp{wieczorek2006})
  dotted. Both bowls bottom inside or at the edge of the admissible
  band.}
  \label{fig:alphasweep}
\end{figure}

\section{Surface-temperature closure against Diviner observations}\label{sec:divinerclosure}
The retrieval fits only the deep sensors, so the diurnal surface
temperature is an out-of-sample check: any admissible column must
also radiate a plausible surface cycle. Figure~\ref{fig:diviner}
closes the loop against the \citet{williams2017diviner} Diviner
Global Cumulative Product (GCP) diurnal composite, the standard
orbital record of lunar surface temperatures, at each site. The Hayne form at the per-site \Kds\
reproduces the composite to $4.1$~K full-cycle RMSE at Apollo~15 and
$6.3$~K at Apollo~17 (mean biases $+2.6$ and $+4.1$~K); the
surface-calibrated \citet{martinez2021} form stays closer to the
surface data, $3.6$~K at both sites with biases $+0.1$ and $-0.8$~K.
That is the mirror image of the deep-sensor ranking, where the
per-site Hayne fit wins.

\begin{figure}[t]
  \centering
  \includegraphics[width=0.96\linewidth]{figures/fig_diviner_closure.pdf}
  \caption[Surface-temperature closure against the Diviner
  composite]{Surface-temperature closure against the
  \citet{williams2017diviner} Diviner GCP diurnal composite (gray
  markers). Teal solid: Hayne form at the per-site retrieved \Kds;
  violet dashed: \citet{martinez2021} at published coefficients. The
  local-solar-time axis is centered on midnight; each panel's inset
  gives the full-cycle RMSE and mean bias, and the nighttime detail
  is enlarged.}
  \label{fig:diviner}
\end{figure}

The Apollo~17 mechanism is direct. The high meter-scale
$\Kds=7.08$~\mWmK\ couples more heat upward through the diurnal skin
than the globally calibrated $K(T,z)$ assumes, raising the modeled
dayside response; the per-site fit accepts that $+4.1$~K mean surface
penalty as the price of matching the deep gradient, while the
surface-calibrated Mart\'inez form remains nearly unbiased at the
surface and pays its penalty at the deep sensors instead ($0.63$
against $0.40$~K). No single parameter choice within the Hayne family
serves both depth regimes at Apollo~17 at once. Both models
nonetheless track the composite through a ${\sim}280$~K swing to
within $3.6$--$6.3$~K, so the insolation, albedo, emissivity, and
shallow $K(T,z)$ chain is captured correctly by both, and the
deep-sensor discrepancy is not a surface-driven artifact.

\section{Robustness of the contrast to the
  \texorpdfstring{$\Kd$--$\Qb$}{Kd-Qb} degeneracy}\label{sec:robustness}

\subsection{Failure of the proportional-degeneracy assumption}
The deep temperature slope depends on the interior heat flow and the
blanket's conductivity together: a stronger flow through a more
conductive blanket can look like a weaker flow through a less
conductive one, so the data alone cannot fully tell the two apart.
This section measures, rather than assumes, how far that ambiguity
extends. \Kd\ and \Qb\ enter the deep gradient together
($d\Tm/dz\simeq\Qb/K$), which suggests the proportional invariance
$(\alpha\Kd,\alpha\Qb)\to$ same profile. Under that assumption a
basal-flux error rescales both sites equally and cancels in the
contrast. Direct re-retrieval disproves the invariance for this
problem: the retrieval also matches absolute sensor temperatures,
whose offset through the low-conductivity skin depends on $\Qb/K$
non-proportionally, and the radiative $K\propto T^3$ coupling breaks
the scaling further. Figure~\ref{fig:qbdeg} shows the realized
$\Kds(\Qb)$ relations, one full re-retrieval per point:
$\Kds(\mathrm{A15})$ \emph{rises} with its flux (from $2.07$ to
$5.44$~\mWmK\ across $\Qb=14$--$26$~\mWm) while $\Kds(\mathrm{A17})$
\emph{falls} (from $10.13$ to $6.80$ across $10$--$18$~\mWm).
Opposite responses of this kind cannot be represented by any
proportional model.

The opposite signs have a physical reading. The retained Apollo~15
band ($80$--$139$~cm) carries a comparatively steep internal
gradient, and the fit tracks that gradient $\Qb/K$, tying \Kds\ to
\Qb\ with a positive sign. The retained Apollo~17 band
($130$--$234$~cm) is deeper and nearly isothermal, so its fit is
controlled by the absolute temperature level, which responds to \Qb\
mainly through the radiative $T^3$ coupling: a higher flux warms the
column, inflates the radiative share of the conductivity, and the
fitted contact term falls to compensate.

\begin{figure}[htbp]
  \centering
  \includegraphics[width=0.95\linewidth]{figures/fig_qb_degeneracy.pdf}
  \caption[The measured $\Kds(\Qb)$ relations]{The measured
  $\Kds(\Qb)$ relation at each site (one full re-retrieval per grid
  point). The two sites respond to a flux revision in opposite
  directions, which is what makes the contrast's sign
  envelope-robust.}
  \label{fig:qbdeg}
\end{figure}

\subsection{Direct mapping of the degeneracy}
The degeneracy is therefore mapped by brute honesty: the contrast
$\Delta\Kds$ over the full published flux envelope
($\Qb^{\mathrm{A15}}\in[14,25]$, $\Qb^{\mathrm{A17}}\in[10,18]$~\mWm),
interpolated from the measured curves of Fig.~\ref{fig:qbdeg} with no
degeneracy model at all (Fig.~\ref{fig:robustness}a). That panel reads
as a map: each point on it asks what the contrast would be if the two
basal fluxes took those values, and the color is the answer, so a
uniformly warm map means the sign of the contrast does not depend on
which fluxes one assumes. The contrast is positive over the entire
measured domain. The most adverse corner,
both fluxes at their upper bounds, compresses it to $+1.5$~\mWmK\
(${\sim}1.6\sigma$); the favorable corner, both at their lower
bounds, expands it to $+8.1$. A $-33\%$ revision of
$\Qb^{\mathrm{A17}}$ alone, the direction the reanalysis debate
points, \emph{raises} the contrast to ${\sim}+5.2$. What the flux
debate does forfeit is the absolute magnitude: the same map moves
$\Kds(\mathrm{A17})$ between ${\sim}6.8$ and ${\sim}10.1$~\mWmK\
across the debated envelope, and at Apollo~17 the lower
\citet{saito2007} flux actually fits the record slightly better (RMSE
$0.29$ against $0.40$~K), so fit quality cannot arbitrate the flux
debate.

\begin{figure}[htbp]
  \centering
  \includegraphics[width=\linewidth]{figures/fig_robustness.pdf}
  \caption[Robustness of the contrast]{Robustness of the contrast.
  \textbf{(a)}~The inter-site contrast
  $\Delta\Kds=\Kds(\mathrm{A17})-\Kds(\mathrm{A15})$ mapped over the
  plane of the two assumed basal fluxes (Apollo~15 flux on the
  horizontal axis, Apollo~17 on the vertical), one full re-retrieval
  per grid point. Warm colors mean Apollo~17 is the more conductive
  site; the contrast is positive across the entire plausible domain.
  The two published flux estimates are marked, Langseth~1976 (black
  dot) and the lower Saito~2007 reanalysis (green square); the dashed
  contours give the contrast's significance in bootstrap standard
  errors ($2\sigma$/$4\sigma$/$6\sigma$); and the teal loops are the
  MCMC posterior for where the true fluxes most likely lie (68\% inner,
  95\% outer), with its median starred. The star is the flux pair the
  data prefer, not a point where two models agree, and it sits below
  Langseth because the record mildly favors lower fluxes. The
  posterior-implied ordering holds with
  $P(\Delta\Kds>0)=99.2$--$99.4\%$. \textbf{(b)},\textbf{(c)}~Joint
  $(\Kd,H)$ misfit surfaces per site, showing the shallow valley along
  which \Kd\ and the compaction depth $H$ trade off.}
  \label{fig:robustness}
\end{figure}

The same direct-map logic extends to the two systematics that rival
\Qb\ in Table~\ref{tab:budget}, and the sites again respond in
opposite directions. Raising the Bond albedo by $0.01$ lowers
$\Kds(\mathrm{A15})$ from $4.60$ to $3.20$ but raises
$\Kds(\mathrm{A17})$ from $7.08$ to $11.81$, so the contrast stays
positive across the full $\pm0.01$ envelope even with the two albedos
perturbed independently (most adverse pairing $+1.3$~\mWmK); the
$\pm30\%$ \Ks\ envelope behaves the same way (most adverse pairing
$+1.6$). These systematics limit each site's absolute \Kds\ without
threatening the sign of the contrast.

\subsection{Bayesian cross-check}\label{sec:bayes}
As an independent check on the direction of the contrast, and not a
second estimate of its magnitude, a Bayesian inversion on the same
directly solved RMSE$(\Kd,\Qb)$ surfaces propagates the
\citet{saito2007}/\citet{nagihara2018} flux priors, the assumed
probable ranges of \Qb, by Markov-chain Monte Carlo (MCMC) sampling
\citep{foreman_mackey2013emcee}. The ordering probability
$P[\Kds(\mathrm{A17})>\Kds(\mathrm{A15})]$ stays at $99.2$--$99.4\%$
for prior widths from $\pm10\%$ to $\pm40\%$ of \Qb, and the contrast
median grows with prior width, from $+3.4$ to $+4.8$~\mWmK. The
ordering is insensitive to the very systematic that dominates the
error budget, precisely because the two sites move oppositely under
it. The marginal MCMC medians, $\Kd(\mathrm{A15})=4.3$ and
$\Kd(\mathrm{A17})=8.3$~\mWmK, differ from the bootstrap point
estimates by construction rather than disagreement: the \Qb\ prior
slides each marginal along the degeneracy ridge, so the two methods
are not independent measurements of the absolute value.

Figure~\ref{fig:postcompare} sets the two uncertainty views side by
side: the flux-agnostic bootstrap, whose contrast interval grazes
zero, and the flux-prior Bayesian posterior, firmly positive.
Together they bracket the honest statement. Conditional on the
published fluxes the magnitude is marginal; once the flux debate
itself is propagated the ordering is decisive, and the absolute
values inherit the full flux uncertainty. The flux posteriors
themselves pull below the published values at both sites (medians
$17.8$ against a $21\pm3.2$ prior at Apollo~15, $13.9$ against
$16\pm2.4$ at Apollo~17), independently favoring the
\citet{saito2007} direction. A firm per-site \Qb\ from a future
landed measurement remains the single most valuable external input
this retrieval could receive.

\begin{figure}[htbp]
  \centering
  \includegraphics[width=\linewidth]{figures/fig_posterior_compare.pdf}
  \caption[Bootstrap vs.\ Bayesian posteriors]{The two uncertainty
  views compared: per-site posterior overlays (left) and the contrast
  posterior (right) from the MCMC cross-check, to be read against the
  bootstrap distributions of Fig.~\ref{fig:boot}.}
  \label{fig:postcompare}
\end{figure}

\section*{Chapter summary}
The per-site retrieval gives $\Kds(\mathrm{A15})=4.60$ and
$\Kds(\mathrm{A17})=7.08$~\mWmK, cutting the deep-sensor misfit of
the global value from $1.09$ to $1.00$~K and from $0.89$ to $0.40$~K
(\S\ref{sec:headline-results}). Model selection prefers the per-site
description decisively at Apollo~17 and is indifferent at Apollo~15
(\S\ref{sec:modelsel}); the fits generalize across thermometer
families and to the withheld deepest sensors (\S\ref{sec:holdout});
an epoch-aligned reduction moves Apollo~17 to $6.16$~\mWmK\ and is
carried as its largest methodological systematic
(\S\ref{sec:commonepoch}); an independent global form misses the same
sensors in the same direction and cannot be rescued by any admissible
density (\S\ref{sec:secondmodel}); and the fitted columns radiate
surface cycles within $3.6$--$6.3$~K of the Diviner composite, so the
deep discrepancy is not a surface artifact
(\S\ref{sec:divinerclosure}). The one thing the boreholes cannot
settle alone, the basal-flux conditionality of the absolute values,
is mapped directly in \S\ref{sec:robustness}: the sign of the
contrast survives the entire published flux envelope, its magnitude
does not.

% ============================================================================
\chapter{Discussion and conclusions}\label{ch:discussion}

\section{Two admissible interpretations of the contrast}\label{sec:tworeadings}
The HFE record does not admit a single deep conductivity that fits
both sites. Under the Hayne form at the published fluxes, the value
that fits Apollo~17 degrades Apollo~15 to a $+0.70$~K bias (RMSE
$1.16$~K), and the value that fits Apollo~15 degrades the Apollo~17
RMSE from $0.40$ to $0.54$~K. What the record constrains tightly at
each borehole is the meter-scale temperature gradient; attributing
the difference in gradients to conductivity rather than to basal flux
rests on the published \Qb\ values. The data therefore admit two
readings, a conductivity contrast at the published fluxes or a flux
contrast at a shared conductivity, and a pooled three-model
comparison cannot separate them: per-site \Kd\ with fixed fluxes,
shared \Kd\ with fitted per-site fluxes, and the strict shared-\Kd\
null all land within $\Delta\mathrm{AICc}\le0.6$ of one another on
the pooled sensor set. The pooled tie is expected rather than
troubling, because pooling dilutes the decisive Apollo~17 evidence
($\Delta\mathrm{AICc}\approx-23$) with the less informative Apollo~15
set ($N=7$). The defensible statement is that the two boreholes
require different meter-scale temperature gradients, decisively at
Apollo~17, and that the conductivity reading of that requirement is
the one consistent with the published basal fluxes.

\section{Physical plausibility of the contrast}\label{sec:physorigin}
A factor-${\sim}1.5$ conductivity contrast between Hadley Rille and
Taurus--Littrow is physically unremarkable. Apollo cores span bulk
porosities of $0.30$--$0.45$ \citep{mitchell1973,carrier1991}, and
the effective-medium scaling $K_{\rm eff}\propto(1-\phi)^2$
\citep{wood2020} yields a factor of ${\sim}1.5$ across that range
alone. Composition points the same way: solid basalt conducts about
$30\%$ better than anorthosite \citep{horai1972}, Apollo~17 sits in a
basaltic embayment, and Apollo~15 straddles a mare/highland margin
\citep{korotev2005,wieczorek2006,crawford2015}. The product of the
two mechanisms, about $1.9$, covers the retrieved ratio without
pushing either to its extreme. The Mart\'inez density diagnostic of
\S\ref{sec:secondmodel} adds a useful negative: the density that form
requires at Apollo~17 ($1877$~kg\,m$^{-3}$) is ordinary, so the
enhanced conductivity is carried by composition and contact
structure, not by anomalous compaction. No site-resolved porosity
contrast has been measured at these boreholes, so this is a
plausibility argument, not a demonstration. What the data establish
is the ordering; what the geology supplies is that a contrast of this
size needs nothing exotic.

Figure~\ref{fig:mechanisms} puts the budget in one frame. Either
mechanism alone nearly reaches the retrieved ratio of $1.54$, and
their product clears it with margin, so the dashed line sits
comfortably inside what ordinary regolith variation supplies. If the
retrieved ratio had landed at $3$ or beyond, this figure would look
very different and the result would demand an exotic explanation; at
$1.54$ it demands none.

\begin{figure}[htbp]
  \centering
  \includegraphics[width=0.72\linewidth]{figures/fig_contrast_mechanisms.pdf}
  \caption[Mechanism budget for the conductivity ratio]{Mechanism
  budget for the retrieved conductivity ratio. Bars: the factor each
  mechanism supplies on its own (porosity across the Apollo-core
  range via $K\propto(1-\phi)^2$; composition from the
  basalt/anorthosite contrast) and their product. Dashed line: the
  retrieved ratio $7.08/4.60=1.54$. Ordinary variation covers the
  contrast without either mechanism at its extreme.}
  \label{fig:mechanisms}
\end{figure}

\section{Relation to prior conductivity estimates}\label{sec:priorwork}
The per-site values land inside the historical spread
(Table~\ref{tab:prior}) and help reconcile its extremes. The
comparison requires one distinction that the literature often blurs:
the \emph{effective} conductivity a sensor at depth measures, with
the radiative term included, against the \emph{contact} asymptote
\Kd\ of the Hayne parameterization. The two differ by the radiative
factor $1+\chi(T/T_{\rm ref})^3\approx2.04$ at the deep-sensor mean
of $255$~K, and values from different classes must not be compared
directly.

\begin{table}[htbp]
  \centering\small
  \caption[Five decades of meter-scale conductivity
  estimates]{Five decades of meter-scale regolith conductivity
  estimates, separated by quantity class: the \emph{effective}
  (total) conductivity a sensor at depth measures, radiative term
  included, versus the \emph{contact} asymptote \Kd\ of the
  Hayne-form parameterization. The classes differ by the radiative
  factor ${\sim}2$ at $255$~K.}
  \label{tab:prior}
  \begin{tabular}{@{}p{0.30\linewidth}lp{0.21\linewidth}l@{}}
    \toprule
    Source & Class & Method & Value (\mWmK) \\
    \midrule
    \citet{cremers1971}, Apollo 12 fines & effective & laboratory, returned fines & $1.2$--$3.5$ \\
    \citet{langseth1972,langseth1976}, original & effective & in-situ HFE, initial & $1$--$3$ \\
    \citet{langseth1976}, revised annual-wave & effective & in-situ HFE, diffusivity & ${\approx}10$ (A15), ${\approx}12.4$ (A17) \\
    \citet{keihm1984}; \citet{grott2010} & effective & HFE reanalyses & upward revision; reassessment \\
    \citet{vasavada2012} & effective, at depth & orbital, Diviner equatorial & ${\approx}7$ \\
    \citet{hayne2017}, global fit & contact \Kd & orbital, Diviner global & $3.4$ \\
    \citet{feng2020}, global fit & contact \Kd & orbital, microwave & $3.8$ \\
    \midrule
    This work (per site) & contact \Kd & in-situ HFE, profile fit & $4.60$ (A15), $7.08$ (A17) \\
    \quad as effective at the sensors & effective & --- & $9.3$ (A15), $14.6$ (A17) \\
    \bottomrule
  \end{tabular}
\end{table}

Expressed as effective conductivities at the sensors, the quantity
the annual-wave reductions actually measure, the retrievals are
$9.3$ and $14.6$~\mWmK. These bracket the revised
\citet{langseth1976} figures of ${\approx}10$ and
${\approx}12.4$~\mWmK: Apollo~15 agrees to ${\sim}7\%$, Apollo~17
sits at the upper edge of the in-situ scatter. The decisive point is
the ordering. Langseth's own reduction already places Apollo~17 above
Apollo~15, from an independent analysis of the same records four
decades earlier, so the ordering retrieved here is not an artifact of
the present forward model. The two analyses differ mainly in the
radiative term and the borestem treatment, which move the absolute
level rather than the sign. Since the historical in-situ reductions
have themselves been critically reassessed \citep{grott2010}, this
work anchors the conductivity in the measured temperature profile
rather than in the annual-wave analysis, and treats the agreement
above as a consistency check, not an input.

\section{Caveats and limitations}\label{sec:caveats}
Eight caveats accompany any use of these results.
\begin{enumerate}
  \item \textbf{Marginal magnitude.} The bootstrap contrast interval
    touches zero ($P_{\rm boot}=0.031$, a bootstrap tail proportion,
    not a null-centered $p$-value). The claim defended is the
    ordering and its robustness, not a $95\%$-significant magnitude.
  \item \textbf{\Qb\ conditionality.} The deep gradient constrains
    $\Qb/\Kd$, so the absolute \Kds\ values are conditional on the
    published fluxes, and the data favor but do not establish a
    conductivity contrast over a uniform conductivity with differing
    fluxes (\S\ref{sec:tworeadings}). The published fluxes also
    derive from the same records, a circularity mitigated by the
    direct map and the flux-agnostic bootstrap.
  \item \textbf{$\chi$ conditionality.} The radiative coefficient is
    held at the global $\chi=2.7$. Re-retrieving at the
    \citet{vasavada2012} normalization $\chi=1.48$ drives the RMSE
    minimum outside the physical sweep range at both sites, so \Kds\
    is conditional on the published $\chi$ in the same sense as on
    \Qb. The contrast is less exposed, since both sites share the
    form.
  \item \textbf{$H$ conditionality.} The joint $(\Kd,H)$ surface is a
    shallow valley (Fig.~\ref{fig:robustness}b,c): the interior
    minima improve on the fixed-$H$ retrieval by only $0.10$~K (A15)
    and $0.09$~K (A17), so $H$ is effectively unconstrained and \Kds\
    is reported at the published $H=6$~cm, with the trade-off carried
    as the conditional $H$ row of Table~\ref{tab:budget}.
  \item \textbf{Epoch mixing.} The stability windows mix calendar
    epochs; the common-1974 reduction moves Apollo~17 by $-0.89$~\mWmK\
    and makes the contrast more marginal ($P_{\rm boot}=0.04$)
    while preserving its sign (\S\ref{sec:commonepoch}). The
    bootstrap also treats sensors as exchangeable; a probe-clustered
    (hierarchical) treatment remains follow-on work.
  \item \textbf{Site--depth ambiguity.} The two sensor sets barely
    overlap in depth ($80$--$139$ against $130$--$234$~cm), so a
    conductivity still increasing below ${\sim}1.5$~m could in
    principle mimic a site contrast. Two internal diagnostics limit
    this reading: within the Apollo~17 set, which itself spans a full
    meter, a depth-constant \Kd\ leaves only $0.40$~K RMSE and
    predicts the withheld $234$~cm sensor to $0.25$~K, where a still
    rising $K(z)$ would curve the profile detectably; and the
    borestem-cut sweep leaves both retrievals unchanged as the
    Apollo~15 set is restricted to its deepest five sensors.
    Compaction below the probed range cannot be excluded with two
    boreholes, but the data do not require it.
  \item \textbf{Topography.} The forcing assumes a flat surface.
    Taurus--Littrow's massifs shadow the Apollo~17 site near the
    terminators and re-radiate toward it at midday; horizon blocking
    removes on the order of $1\%$ of the lunation-mean insolation at
    a valley-floor site, the same order as the albedo envelope,
    partially offset by terrain infrared and reflected sunlight.
    Terrain effects plausibly contribute to the $+4.1$~K Apollo~17
    surface bias of \S\ref{sec:divinerclosure}; a terrain-resolved
    forcing is the natural refinement.
  \item \textbf{Two points, not a map.} Two boreholes cannot
    constrain global variability. They demonstrate that at least this
    much site-to-site variation exists.
\end{enumerate}

\section{Implications for subsurface-sounding missions}\label{sec:implications}
Any mission that senses below the diurnal skin needs a meter-scale
$T(z)$ boundary condition that orbital conductivity models cannot
supply on a per-site basis. The retrieved per-site columns provide
that constraint from the only in-situ data available, and the
immediate consumer is TSUKIMI \citep{tsukimi2022}: its terahertz
forward model ingests exactly the sub-skin temperature column this
retrieval anchors. Figure~\ref{fig:tsukimichain} draws the chain.
The link this thesis supplies is the last one: everything upstream of
it exists or is under construction on the mission side, and the
$T(z)$ hand-off is the point where in-situ ground truth enters an
otherwise orbital pipeline.

Two limits on the hand-off should be stated as plainly as the
hand-off itself. The anchors are two nearside sites at $26^{\circ}$
and $20^{\circ}$N; carrying their \Kds\ values to the polar latitudes
TSUKIMI will map rides on the spatial assumptions of the very global
form this thesis has shown to miss both in-situ anchors. What
transfers is therefore the demonstration that site-resolved values
are needed, the certified machinery to retrieve them wherever future
landed data exist, and the two anchor columns themselves as
validation points. The deliverable is concrete: the pipeline of
Appendix~\ref{app:repro} regenerates each site's converged $T(z)$
column from the archived configuration, in a form any forward model
can ingest.

\begin{figure}[htbp]
  \centering
  \includegraphics[width=\linewidth]{figures/fig_tsukimi_chain.pdf}
  \caption[Where this work enters the TSUKIMI pipeline]{Where this
  work enters a subsurface-sounding pipeline. Terahertz emission
  originates below the diurnal skin, a layer no orbital infrared
  calibration constrains, so the mission's forward model needs a
  meter-scale $T(z)$ column; the per-site \Kds\ retrieved here
  anchors that column in the only in-situ data.}
  \label{fig:tsukimichain}
\end{figure}

The error budget also tells such missions where
the leverage is. The dominant terms are the basal flux and the
albedo, not the sensors, so a single firm per-site \Qb\ from a future
landed heat-flow measurement would tighten the retrieval more than
any amount of reprocessing; the solver already accepts a per-site
flux as a boundary condition, so no change of method would be
needed.

\section{Conclusions}\label{ch:conclusion}
This thesis set out to test one number, the globally averaged deep
conductivity $\Kd=3.4$~\mWmK, against the only in-situ meter-scale
temperatures the Moon affords, and to build the machinery that makes
the test affordable.

The three research questions of Chapter~\ref{ch:intro} receive direct
answers. \emph{What deep conductivity does each borehole require?}
$\Kds(\mathrm{A15})=4.60^{+0.86}_{-0.21}$ and
$\Kds(\mathrm{A17})=7.08^{+0.53}_{-0.52}$~\mWmK, with total $1\sigma$
budgets of $\pm1.9$ and $\pm3.9$, conditional on the published basal
fluxes and the Hayne radiative normalization. \emph{Do the two sites
require the same value?} No at the point estimates, and the ordering
$\mathrm{A17}>\mathrm{A15}$ survives the entire published flux
envelope; the magnitude of the difference, however, is marginal at the
$95\%$ level and is reported as such. \emph{Can the test be made
computationally routine?} Yes: the flux-anchored solver cuts the cost
of one periodic-equilibrium solve by about a factor of $2500$, which
turned the bootstrap, the degeneracy map, and the sensitivity sweeps
from a prohibitive proposition into a routine one. Four results carry
these answers.

\begin{itemize}
  \item The Apollo 15 and 17 boreholes favor deep conductivities of
    $4.60^{+0.86}_{-0.21}$ and $7.08^{+0.53}_{-0.52}$~\mWmK. The
    global $3.4$~\mWmK\ falls outside both $95\%$ bootstrap intervals
    at the published inputs, and remains about $1\sigma$ away once
    the full systematic budget is folded in. On the retrieval's own
    terms, meter-scale lunar thermal modeling should carry
    site-resolved rather than global \Kd.
  \item The inter-site contrast (median $+2.31$~\mWmK) is marginal in
    magnitude and robust in direction. Its sign survives the entire
    published basal-flux envelope, with an MCMC ordering probability
    of $99.2$--$99.4\%$ regardless of prior width, because the two
    sites respond to flux revisions in opposite directions, a fact
    found by mapping the degeneracy directly instead of assuming its
    proportional form. The ordering is corroborated by the
    independent annual-wave analysis of the same records by
    \citet{langseth1976}, which placed Apollo~17 above Apollo~15 four
    decades earlier.
  \item The flux-anchored solver makes exhaustive uncertainty
    analysis of borehole-constrained thermal models routine
    (about ${\sim}2500\times$ faster per solve than brute force), and
    is certified against brute force, a time-step ladder on the
    retrieved quantity, and an independent C++ implementation to
    $10^{-12}$~K.
  \item For TSUKIMI and successor subsurface sounders, the retrieved
    per-site columns supply a meter-scale temperature constraint
    orbital instruments cannot measure. In the other direction, a
    firm per-site \Qb\ from a landed measurement is the one external
    input that would most tighten this retrieval.
\end{itemize}

Two themes connect the results. The method made the measurement
defensible: without the thousandfold speed-up, neither the 1500-draw
bootstrap nor the point-by-point degeneracy map would have been
affordable, and the marginal contrast would rest on a single
deterministic fit. The affordability in turn made the honesty cheap.
Because every number could be re-run under every perturbed input, the
caveats of \S\ref{sec:caveats} are quantified rather than asserted,
and the one claim defended without qualification, the
$\mathrm{A17}>\mathrm{A15}$ ordering, is the one the data support
across the published flux envelope.

\subsection{Future work}\label{sec:future}
Four extensions follow directly.

A landed heat-flow measurement at either site would fix the input
that limits the contrast. The solver takes a per-site flux as a
boundary condition, so a firm \Qb\ would turn the marginal magnitude
into a decided one with no methodological change.

The same sweep can be run under the \citet{martinez2021} form, with a
per-site recalibration of its contact/radiative coefficients against
the HFE record. The density diagnostic of \S\ref{sec:secondmodel}
shows where its published form saturates; comparing retrieved deep
values across the two forms would separate regolith properties from
parameterization artifacts.

The radiative coefficient $\chi$ and \Kd\ are partly degenerate in a
fit to lunation means. A joint retrieval constrained by daytime
brightness temperatures or an independent density profile would close
the largest conditionality this thesis leaves open.

Finally, the solver's speed makes a spatial extension practical, and
this is the direction I intend to pursue. The two-borehole retrieval
is a point measurement; the same forward model, driven pixel by
pixel, would produce meter-scale $T(z)$ maps rather than two columns.
Three steps lead there. The first is a terrain-resolved surface
forcing: at valley-floor and polar sites the flat-surface insolation
of \S\ref{sec:neglected} breaks down, so horizon shadowing and
terrain re-radiation computed from a lunar digital elevation model
must replace it, most critically for the permanently shadowed regions
near the poles. The second is a global, or at least polar, map version
of the pipeline, forward-modeling the subsurface column at the
spatial resolution a mission such as TSUKIMI requires. The third
couples that thermal map to volatile transport: with a per-pixel
$T(z)$ in hand, sublimation and cold-trapping rates yield an
ice-survivability map, the depth at which water ice is stable against
loss over geological time. That chain, from in-situ-anchored
conductivity to a lunar water-cycle product, is the longer-term aim
this thesis lays the groundwork for.

% ============================================================================
\appendix

\chapter{Derivation of the closure ODE}\label{app:derivation}

This appendix derives Eq.~\eqref{eq:closure} from the heat equation
in five moves, none skipped.

\section*{Move 1: average over one lunation}
Define the cycle mean $\langle\,\cdot\,\rangle \equiv
\frac{1}{P}\int_0^P (\cdot)\,dt$ over one lunation $P$. In periodic
steady state the column repeats every cycle, $T(z,t+P)=T(z,t)$. Apply
$\langle\,\cdot\,\rangle$ to both sides of Eq.~\eqref{eq:heat}.

\section*{Move 2: the storage term vanishes}
Write the volumetric enthalpy $h(T)=\int\rho\,c_p\,dT$, so that
$\rho c_p\,\partial_t T=\partial_t h$. Then
\begin{equation}
  \left\langle \rho c_p\,\frac{\partial T}{\partial t}\right\rangle
  = \frac{1}{P}\int_0^P \frac{\partial h}{\partial t}\,dt
  = \frac{h(P)-h(0)}{P} = 0,
\end{equation}
because $h$ returns to its starting value after a full cycle. In
steady state nothing accumulates heat on average: time drops out of
the problem. This move is legal only in periodic steady state, which
is why Step~A must first drive the skin to that state before Step~B
is invoked.

\section*{Move 3: the depth derivative passes through the average}
On the right-hand side, $z$ and $t$ are independent variables, so the
average commutes with $\partial_z$:
\begin{equation}
  0 = \left\langle \frac{\partial}{\partial z}
      \left(K\,\frac{\partial T}{\partial z}\right)\right\rangle
    = \frac{\partial}{\partial z}
      \left\langle K\,\frac{\partial T}{\partial z}\right\rangle
  \quad\Longrightarrow\quad
  \left\langle K\,\frac{\partial T}{\partial z}\right\rangle
  = \text{constant in } z.
\end{equation}
No $t$ remains: the averaged equation is an equation in depth alone.

\section*{Move 4: pin the constant with the bottom boundary}
The basal boundary condition fixes the constant: at $z_{\max}$ the
conducted flux is \Qb, hence
\begin{equation}
  \left\langle K\,\frac{\partial T}{\partial z}\right\rangle = \Qb
  \qquad\text{at every depth.}
\end{equation}
The cycle-mean column carries the geothermal flux all the way up.

\section*{Move 5: peel the average off the product}
The mean of a product is not the product of the means; name the gap
\begin{equation}
  \urect(z) \;\equiv\;
  \left\langle K\,\frac{\partial T}{\partial z}\right\rangle
  - K(\Tm, z)\,\frac{d\Tm}{dz},
  \label{eq:urect}
\end{equation}
the \emph{rectified flux}, the covariance of the oscillating
conductivity and the oscillating gradient, nonzero only where the
diurnal wave lives. Substituting Move~4 into Eq.~\eqref{eq:urect} and
solving for the mean gradient gives the closure,
\begin{equation}
  \frac{d\Tm}{dz} = \frac{\Qb - \urect(z)}{K(\Tm, z)},
\end{equation}
a first-order ODE in depth alone. Above the anchor, \urect\ is
evaluated from the Step~A skin solution; below the anchor it decays
rapidly (at $z_0=0.55$~m the diurnal amplitude is ${\sim}0.15$~K and
$\urect<1\%$ of \Qb), so the deep integration is effectively
$d\Tm/dz=\Qb/K$.

\section*{Step B in practice}
Equation~\eqref{eq:closure} is integrated downward from the anchored
temperature $\Tm(z_0)$ by second-order Runge--Kutta on the same
geometric grid used by Step~A, with $K$ evaluated at the current mean
temperature at each half-step. Because the ODE \emph{is} the
steady-state equation, Step~B is exact below the anchor given a
correct anchor temperature; the outer loop's only job is to make the
anchor temperature correct.

% ============================================================================
\chapter{Numerical scheme and certification}\label{app:numerics}

\section{Spatial discretization}
\emph{Grid.} The column $[0,\,5~\mathrm{m}]$ is divided into 69
finite-volume cells whose widths grow geometrically ($\Delta
z_0=2$~mm, growth $8\%$): millimeter resolution where the
${\sim}280$~K diurnal swing lives, centimeter-to-decimeter cells in
the smooth deep column. Temperatures live at cell centers; fluxes
cross cell faces. Integrating the heat equation over a cell and
telescoping the interior fluxes conserves energy to machine
precision.

\emph{Face conductivity.} Heat crossing the face between cells $i$
and $i{+}1$ passes through two half-cells in series, so thermal
resistances add and the face value is the \emph{harmonic} mean,
\begin{equation}
  K_{i+1/2} = \frac{2\,K_i\,K_{i+1}}{K_i+K_{i+1}},
\end{equation}
the physically correct choice where $K$ jumps steeply through the
compaction zone (an arithmetic mean would let a high-$K$ cell
short-circuit its low-$K$ neighbor).

\section{Time integration}
\emph{Scheme.} Crank--Nicolson (the $\theta=\tfrac12$ scheme):
unconditionally stable (the amplification factor per diffusive mode
is $\xi=(1-\beta/2)/(1+\beta/2)$ with $\beta\ge0$, so $|\xi|\le1$ for
every mode and any $\Delta t$) and second-order accurate. Explicit
Euler at the $2$~mm surface cell would demand $\Delta t\lesssim200$~s
for stability alone; CN frees the step to be chosen by accuracy,
certified below at $\Delta t=1800$~s.

\emph{Nonlinearity.} Properties depend on $T$, so each step performs
midpoint-Picard corrector sweeps: properties are re-evaluated at
$\tfrac12(T^{\,n}+T^{\,n+1})$ and the tridiagonal solve repeated. The
production setting of two sweeps was fixed by a measured scan of the
periodic attractor's residual drift per time-step halving:
$-8.6/-4.2/-2.2/-2.0$~mK at $1/2/4/8$ sweeps.

\emph{Surface.} The nonlinear radiative surface balance is solved by
Newton iteration each step, re-anchored at the implicit surface
temperature.

\emph{Linear algebra.} Each step assembles one tridiagonal system
$A\,T^{\,n+1}=d$: every interior row is the CN stencil; the surface
row is a Dirichlet ghost (the Newton-solved surface temperature moves
to the right-hand side); the basal row is a Neumann ghost
($b_{N-1}\to b_{N-1}-\tfrac12\alpha$,
$d_{N-1}\to d_{N-1}+\Delta t\,\Qb/C$). The Thomas algorithm solves it
in $O(N)$.

\section{Certification of the reported quantity}
The convergence criterion throughout is the stability of the
\emph{retrieved} \Kds, the quantity actually reported, never a proxy
such as a fixed cycle count or the flux-closure residual alone. Three
certification campaigns are on record.

\paragraph{The fixed-count trap.} An early configuration declared
convergence after a fixed number of spin-up lunations. The skin
settles in ${\sim}10$ lunations while the deep column is still tens
of kelvin away, so a cycle-to-cycle test on the skin passes long
before the profile has arrived (Fig.~\ref{fig:f1bug}). The cure is
the anchor-drift tolerance of \S\ref{sec:fixedpoint}.

\begin{figure}[htbp]
  \centering
  \includegraphics[width=0.85\linewidth]{figures/fig_book_f1_bug.pdf}
  \caption[The fixed-count convergence trap]{The fixed-count trap:
  the shallow column equilibrates quickly while the deep column is
  still relaxing, so a cycle-to-cycle test on the skin passes
  prematurely.}
  \label{fig:f1bug}
\end{figure}

\paragraph{Inner-loop resolution and the truncated Step A.} The
retrieved \Kds\ drifts monotonically with the inner spin-up length
$n_{\rm inner}$ when Step~A time-steps the full column, because the
slow deep relaxation leaks into the anchor. The structural fix is the
truncation of \S\ref{sec:stepA}; with it, Apollo~15 (the
slower-converging site, lower $K$ hence lower diffusivity) reads
$4.600/4.605/4.606$~\mWmK\ at $n_{\rm inner}=64/96/128$, and the
truncated solve was validated unbiased against the full-column solve
($4.606$ against $4.608$ at $n=128$). Production uses
$n_{\rm inner}=96$.

\paragraph{The time-step ladder and the wrap step.} Halving the time
step $1800\to900\to450$~s moves \Kds(A17) by $0.018$ then
$0.007$~\mWmK\ ($7.0795/7.0975/7.1045$; Fig.~\ref{fig:dtladder}),
inside the $0.05$ acceptance band and consistent with second order.
This ladder is itself the instrument that caught the deepest defect
of the project: the original march looped over a periodic grid
spanning $[0,\,P-\Delta t]$ and relabeled the $P-\Delta t$ state as
$t=0$, a phase slip of exactly $\Delta t$ per lunation that made the
periodic attractor first-order in $\Delta t$ (${+}61$~mK per halving)
regardless of the second-order stepper inside. The fix is a wrap step
closing each cycle. As a cross-check, Richardson extrapolation of the
\emph{defective} scheme's ladder converges to $7.104$~\mWmK, the
fixed scheme's finest-rung value to four digits: two different
discretizations, one continuum limit.

\begin{figure}[htbp]
  \centering
  \includegraphics[width=0.85\linewidth]{figures/fig_dt_ladder.pdf}
  \caption[The $\Delta t$ certification ladder]{The time-step
  certification ladder for \Kds(A17) under the wrap-corrected
  scheme.}
  \label{fig:dtladder}
\end{figure}

\paragraph{Independent implementation.} A faithful C++ port of the
solver (same grid, scheme, and wrap step) agrees with the Python
implementation to $1.1\times10^{-12}$~K over full retrievals, and runs
a constant factor of about $1.4\times$ faster per lunation than the
compiled Python kernel, confirming the remaining gap is a language
difference and no algorithmic slack remains.

\section{The audit trail}
Table~\ref{tab:audittrail} records the three input errors found
during the project's audits and their effect on the Apollo~17
headline value. The record is preserved because each is a generic
failure mode of borehole-scale thermal modeling: duplicated constants
drift, convergence criteria certify proxies, and time marches acquire
off-by-one phase defects that only an end-to-end ladder on the
reported quantity exposes.

\begin{table}[h!]
  \centering\small
  \caption{The audit trail: input errors found, their fixes, and the
  resulting moves in \Kds(A17) (\mWmK).}
  \label{tab:audittrail}
  \begin{tabular}{llll}
    \toprule
    Error & Symptom & Fix & \Kds(A17) \\
    \midrule
    stale \Qb\ (15 vs.\ 16~\mWm) & none (ran silently) &
      single source of truth & $8.12\to7.51$ \\
    $n_{\rm inner}$ under-resolved & \Kds\ drifts with $n_{\rm inner}$ &
      truncated Step A + plateau test & $7.51\to7.16$ \\
    missing wrap step & attractor first-order in $\Delta t$ &
      wrap step + $\Delta t$ ladder & $7.16\to7.08$ \\
    \bottomrule
  \end{tabular}
\end{table}

Together the three rows moved the headline number by a full milliwatt,
comparable to the statistical interval itself, which is the practical
argument for auditing inputs with the same rigor as physics. None of
the three produced a crash, a warning, or an implausible profile;
each was caught only because a certification target was defined on
the reported quantity and re-derived from scratch.
Figure~\ref{fig:auditwf} shows the same three corrections as a
waterfall in \Kds(A17).

\begin{figure}[htbp]
  \centering
  \includegraphics[width=0.62\linewidth]{figures/fig_audit_waterfall.pdf}
  \caption[The audit trail as a waterfall]{The three audited input
  errors as a waterfall in \Kds(A17): the stale basal flux, the
  under-resolved inner loop, and the missing wrap step together moved
  the headline value by a full milliwatt before certification
  (green). Numbers from the audit trail of
  Table~\ref{tab:audittrail}.}
  \label{fig:auditwf}
\end{figure}

\section{Regression guards}
The certification campaigns above are one-time demonstrations; the
test suite keeps their conclusions true as the code evolves. Each of
the failure modes in Table~\ref{tab:audittrail}, and each physics
invariant the solver depends on, has a dedicated test among the 49 in
the repository: the basal boundary must be the \Qb\ Neumann ghost and
not a zero-flux copy, face conductivities must be harmonic means, the
closure must carry the $\Qb-\urect$ sign, the time march must include
the wrap step, the truncated and full-column Step~A constructions
must agree on \Kds, and the compiled kernel must reproduce the
interpreted solver. A change that silently reintroduces any audited
defect fails the suite rather than the next audit, which is the
difference between a certified result and a result that happened to
be checked once.

% ============================================================================
\chapter{Benchmark environment and measured timings}\label{app:timings}

The speed-up factors of \S\ref{sec:cost} are ratios of two methods run
on the same machine, so they are independent of the hardware and are
the right way to state the method's cost. The absolute wall-clock
times behind those ratios are not machine-independent, and are given
here only for a reader's concrete sense of scale. They should be read
as illustrative of one setup, not as a property of the method: on
faster or slower hardware every entry in Table~\ref{tab:timings} would
change, while the factors of Table~\ref{tab:speed} would not.

All timings were measured on the author's MacBook Neo (6-core Apple
A18 Pro, 8~GB), as the median of three runs with the compiled
(just-in-time) kernel warmed beforehand so that compilation is
excluded from the timed region. The
brute-force cost is the measured per-lunation march time multiplied by
the ${\sim}3000$ lunations needed to reach ${\sim}0.1$~K of the
periodic steady state; the sweep entries scale one solve by the
$61$ grid points of the two-site retrieval.

\begin{table}[htbp]
  \centering\small
  \caption[Measured wall-clock timings]{Measured wall-clock timings on
  the author's MacBook Neo (6-core Apple A18 Pro, 8~GB; median of three
  runs, kernel warmed). \emph{These values are hardware-specific}; only
  the ratios of Table~\ref{tab:speed} transfer between machines.}
  \label{tab:timings}
  \begin{tabular}{lccc}
    \toprule
    & brute force & flux-anchored & flux-anchored \\
    & (interpreted) & (interpreted) & (compiled kernel) \\
    \midrule
    one solve         & ${\sim}26$~min & $74.7$~s & $0.64$~s \\
    61-solve sweep    & ${\sim}27$~h (proj.) & $76$~min & $39$~s \\
    per lunation      & $523$~ms & \multicolumn{2}{c}{$3.9$~ms (kernel)} \\
    \bottomrule
  \end{tabular}
\end{table}

The independent C++ port, run on the same machine, marches the solver
at $2.8$~ms per lunation against the compiled Python kernel's
$3.9$~ms, a constant factor of about $1.4\times$ that is a language
difference and not an algorithmic one; the two agree to
${\sim}10^{-12}$~K over full retrievals
(\S\ref{sec:certsummary}). Because the bootstrap and the degeneracy
map reuse cached sweep profiles, they add only a small constant on top
of the sweep cost and are not tabulated separately.

% ============================================================================
\chapter{Reproducibility}\label{app:repro}
All results derive from the pipeline in the \texttt{Lunar-HFE}
repository, in three stages:
\begin{enumerate}
  \item \texttt{pipeline/compute/retrieve\_kd.py}: the sweep,
    bootstrap, joint $(\Kd,H)$ grids, and holdout checks, writing
    \texttt{results/kd\_retrieval\_results.json};
  \item the auxiliary computes (degeneracy map, Bayesian cross-check,
    error budget, sensitivities, benchmarks, discretization checks),
    each writing its own \texttt{results/*.json};
  \item \texttt{pipeline/make\_all\_figures.py}: every figure in this
    thesis regenerates from those artifacts; no figure computes
    physics directly.
\end{enumerate}
The test suite (49 tests) covers the physics invariants (boundary
conditions, harmonic faces, closure sign, wrap step, fast-path
equivalence) and must be green before any result is quoted. The HFE
records are the \citet{nagihara2018} restorations; ephemerides follow
\citet{acton1996,acton2018,mazarico2011}.

From a fresh clone, the complete reproduction is:
\begin{quote}\ttfamily
git clone https://github.com/rp3gregorio/Lunar-HFE.git\\
cd Lunar-HFE\\
python3 -m venv .venv \&\& source .venv/bin/activate\\
make install\phantom{xxxxxxxx}\# editable install + dev deps\\
python code/pipeline/fetch\_diviner.py\phantom{xx}\# one-time
$\sim$310 MB\\
make test\phantom{xxxxxxxxxxx}\# 49 tests\\
make retrieve\phantom{xxxxxxx}\# headline $K_d^{*}$ + bootstrap\\
make aux\phantom{xxxxxxxxxxxx}\# sensitivities, MCMC, budget\\
make figures \&\& make paper
\end{quote}
\noindent The canonical retrieval JSON ships with the repository, so
every quoted number can be verified against the committed artifact
without re-running the full sweep; the independent C++
implementation is built with a single \texttt{clang++ -O3} line and
compared by \texttt{benchmark\_cpp.py} (agreement
${\sim}10^{-12}$~K). The complete code, data, and pipeline are openly
available at \url{https://github.com/rp3gregorio/Lunar-HFE}, which is
the archive of record for this thesis.

% ============================================================================
\bibliographystyle{plainnat}
\bibliography{references}

\end{document}
